\documentclass[reqno]{amsart}
\usepackage{a4wide}
\usepackage{hyperref}

\AtBeginDocument{{\noindent\small
\emph{Electronic Journal of Differential Equations},
Vol. 2026 (2026), No. 10, pp. 1--22.\newline
ISSN: 1072-6691. URL: https://ejde.math.txstate.edu, DOI: 10.58997/ejde.2026.10}
\thanks{\copyright 2026. This work is licensed under a CC BY 4.0 license.}
\vspace{8mm}}

\begin{document}
\title[\hfilneg EJDE-2026/10\hfil 
2D chemotaxis-fluid system with double chemical signals]
{Existence and boundedness of the classical global solution to 2D 
chemotaxis-fluid system with double chemical signals}

\author[S. Liu, D. Li, J. Zheng \hfil EJDE-2026/10\hfilneg]
{Shengquan Liu, Dongpu Li, Jiashan Zheng}

\address{Shengquan Liu \newline
School of Mathematics and Statistics,
Liaoning University,
Shenyang 110036,  China}
\email{shquanliu@163.com}

\address{Dongpu Li \newline
School of Mathematics and Statistics,
Liaoning University,
Shenyang 110036,  China}
\email{15142132001@163.com}

\address{Jiashan Zheng (corresponding author) \newline
School of Mathematics and Information Sciences,
Yantai University,
Yantai 264005,  China}
\email{zhengjiashan2008@163.com}

\thanks{Submitted October 21, 2025. Published February 10, 2026.}
\subjclass[2020]{35K20, 35K55, 92C17}
\keywords{Keller-Segel system; boundedness; global existence; sub-logistic source}

\begin{abstract}
This article concerns the  chemotaxis-fluid system with double chemical signals,
\begin{gather*}
n_t + u \cdot \nabla n=\Delta n-\chi \nabla \cdot (n\nabla c)+\xi \nabla \cdot (n\nabla v) + f(n), \quad x \in \Omega, \; t > 0, \\
c_t + u \cdot \nabla c=\Delta c-nc, \quad x \in \Omega, \;t > 0, \\
v_t + u \cdot \nabla v= \Delta v-v+n, \quad x \in \Omega, \; t > 0, \\
u_t =\Delta u+\nabla P+n\nabla \phi,  \quad x \in \Omega, \; t > 0, \\
\nabla \cdot u = 0,  \quad  x \in \Omega, \; t > 0,
\end{gather*}
in a smooth bounded domain $\Omega\subset \mathbb{R}^2$ with
no-flux/no-flux/no-flux/no-slip boundary conditions.  Here $f(n)$ is a given
sub-logistic source function satisfying $f(s)\geq \zeta s$ for small $s\geq0$,
where $\zeta\in \mathbb{R}$, and the growth condition
$$
\liminf_{s \to \infty} \big\{-f(s) \frac{\ln s}{s^2}\big\} = \mu\in (0, \infty].
$$
We obtain the existence and uniform-in-time boundedness of classical global solutions
to the corresponding initial-boundary value problem. We assume that
 the initial data and the physical coefficients satisfy
$$
\left(2 + \chi^2 + \xi^2 + 4\|c_0\|_{L^\infty(\Omega)}^2
+ 4 C_{\mathrm{p}}^2\|c_0\|_{L^\infty(\Omega)}^2
\|\nabla \phi\|_{L^\infty(\Omega)}^2 -\mu\right)^+
< \frac{1}{2 C_{\mathrm{GN}}^4 M},
$$
where $M$ is a constant,
\begin{equation*}
M := \| n_0 \|_{L^1(\Omega)} + |\Omega| \inf_{\eta > 0} \frac{\sup \{ f(s) + \eta s :
s > 0 \}}{\eta},
\end{equation*}
$C_{\mathrm{GN}} > 0$ is the Gagliardo-Nirenberg constant, and
$C_\mathrm{p}$ the Poincar\'e constant. This result reveals that the sub-logistic
source plays a crucial role in preventing blow-up phenomena within the chemotaxis-fluid
system with double chemical signals.
\end{abstract}

\maketitle
\numberwithin{equation}{section}
\newtheorem{theorem}{Theorem}[section]
\newtheorem{lemma}[theorem]{Lemma}
\newtheorem{remark}[theorem]{Remark}
\allowdisplaybreaks

\section{Introduction}

This article concerns the initial-boundary value problem of the chemotaxis-fluid system with double chemical signals and sub-logistic source, which is formulated as follows:
\begin{equation}\label{eq1.1}
\begin{gathered}
n_t + u \cdot \nabla n = \Delta n - \chi \nabla \cdot (n\nabla c)
+\xi \nabla \cdot (n\nabla v) + f(n), \quad x \in \Omega, t > 0, \\
c_t + u \cdot \nabla c = \Delta c - nc, \quad x \in \Omega, t > 0, \\
v_t + u \cdot \nabla v= \Delta v - v + n, \quad
 x \in \Omega, t > 0, \\
u_t  = \Delta u + \nabla P + n\nabla \phi, \quad \nabla \cdot u = 0, \quad
 x \in \Omega, t > 0, \\
\frac{\partial n}{\partial \nu}=\frac{\partial c}{\partial \nu}
= \frac{\partial v}{\partial \nu}= 0, \ u= 0, \quad x \in \partial\Omega, t >
0, \\
n(x, 0) = n_0(x), \; c(x, 0) = c_0(x), \; v(x, 0) = v_0(x), \; u(x, 0) = u_0(x), 
\quad  x \in \Omega.
\end{gathered}
\end{equation}
Here, $\Omega\subset \mathbb{R}^2$ is a bounded domain with smooth boundary 
$\partial\Omega$, and $\nu$ is the unit outward normal vector to $\partial\Omega$. 
The unknown function $n(x,t)$ denotes the density of cells, $c(x,t)$ the concentration 
of oxygen, and $v(x,t)$ the
concentration of the chemical signal secreted by the cells. 
$u(x,t)$ and $P(x,t)$ stand for  the fluid velocity field and its associated pressure,
respectively. In addition, $\phi$ is a given potential function satisfying
\begin{equation}\label{eq1.2}
\phi \in  C^{1 +\kappa}(\Omega) \quad\text{for some } \kappa> 0,
\end{equation}
and $f(n)$ is a sub-logistic source depending $n$, which satisfies the condition 
\eqref{10.20-4} below. The physical constants
$\chi$ and $\xi$ quantify the interaction strengths between the cell and its 
environment:
$\chi>0$ denotes oxygen as an attractor to the cell, while
$\chi<0$ signifies repulsion;
$\xi>0$ indicates that chemical signals secreted by the cell act as repellents, whereas
$\xi<0$ defines them as attractors.

In this article, we aim to prove the global boundedness of the classical solutions 
to \eqref{eq1.1} subject to the initial data $(n_0, c_0, v_0, u_0)$
satisfying
\begin{equation}\label{eq1.3}
\begin{gathered}
0\leq n_0 \in C(\bar{\Omega}), \quad
0\leq c_0 , \quad v_0 \in W^{1, \infty}(\Omega), \\
u_0 \in D(\mathcal{A}^\alpha) \quad \text{for some } 
\alpha \in (\frac{1}{2}, 1),
\end{gathered}
\end{equation}
where $\mathcal{A}$ is the Stokes operator \(\mathcal{A} := -\mathcal{P}\Delta\),
with domain
\[
D(\mathcal{A}) := W^{2,2}(\Omega) \cap W_0^{1,2}(\Omega) \cap L_\sigma^2(\Omega).
\]
Here, \(L_\sigma^2(\Omega)\) represents the space of square-integrable divergence-free 
vector fields, 
\[
L_\sigma^2(\Omega) := \{\varphi \in L^2(\Omega) \mid \nabla \cdot \varphi 
= 0\},
\]
and \(\mathcal{P}\) stands for the Helmholtz projection from 
\(L^2(\Omega)\) onto \(L_\sigma^2(\Omega)\) (see \cite{Sohr-2001}).

To state our main theorem, we first recall some key developments related to
system \eqref{eq1.1} and its variations.

\subsection*{Keller-Segel-type chemotaxis system}  
Chemotaxis is a biological phenomenon characterized by the oriented movement of 
cells in response to chemical signals. The movement towards a higher concentration of 
the chemical substance is termed positive chemotaxis, and the movement towards regions 
of lower chemical concentration is called negative chemotaxis.
The interesting phenomena in diverse biological processes has attracted a lot of 
attention from the biological and mathematical communities  (see \cite{Hillen}). 
Among these chemotaxis models, the general minimal chemotaxis system
 \begin{equation}\label{eq:7.14-1}
\begin{gathered}
n_t = \Delta n - \xi \nabla \cdot (n \nabla v)+f(n), \quad
 x \in \Omega\subset \mathbb{R}^N,\; t > 0, \\
v_t = \Delta v  - v+ n , \quad x \in \Omega\subset \mathbb{R}^N,\; t > 0,
\end{gathered}
\end{equation}
has been extensively studied. 
For $f(n)= 0$, Keller and Segel  \cite{Keller-1970,Keller-1970-2} first proposed 
system \eqref{eq:7.14-1} in the 1970s. 
Since their classical work, the Keller-Segel model \eqref{eq:7.14-1} and its variants 
have been widely investigated by numerous researchers, with substantial advances made 
in understanding the global existence and uniform boundedness of solutions.
When the space dimension $ N = 1 $, all the solutions are global and bounded 
(see \cite{Osaki-2001}). However, if $ N=2$ or $ N \geq 3 $, the solution to the system 
\eqref{eq:7.14-1} exhibits a remarkable feature of finite/infinite time blow-up, 
which arises from the competitive interplay between the diffusion term and the 
aggregation term $-\nabla \cdot (n \nabla v)$. For $ N = 2 $, there exists a threshold: 
if the mass $\int_\Omega n_0$ is less than this threshold, the solution is global and 
bounded (see \cite{G-3, G-4}); conversely, if $\int_\Omega n_0$ exceeds this threshold, 
the system may admit a solution that blows up in finite/infinite time (see \cite{G-5, G-6}). For $N \geq 3$, if $\int_\Omega n_0 > 0$, the solutions to the system \eqref{eq:7.14-1} blow up in finite time for some given radially symmetric positive initial data (see \cite{Herrero-1997, Dolbeault-2004, Nagai}). For more relevant results, we may refer to \cite{G-8, G-9, G-10} and the references therein.
On the other hand, the blow-up phenomena can be ruled out if a suitable logistic source 
$f(n)$ is introduced. Indeed, taking $f(n)=an-bn^2$ with $a, b \in
\mathbb{R}$ in the system \eqref{eq:7.14-1}, all the solutions are global and uniformly 
bounded for arbitrary \( b > 0 \) on \( N = 1, 2 \) (see
\cite{Xiang2018-C,Osaki-2002}). 
For \( N \geq 3 \), Winkler \cite{Winkler2010} found that the Neumann boundary value 
problem of the system \eqref{eq:7.14-1} admits a global bounded classical
solution provided that \( b>b_0\) for some sufficiently large constant $b_0>0$. 
Later, Xiang \cite{Xiang2018-J} improved the result in \cite{Winkler2010} by giving an
explicit expression of the low bound as \( b_0=9\chi/(\sqrt{10} - 2) \) 
in some non-convex domains. Recently, in the case of \( f(n) = an -b n^\alpha\),  
Winkler \cite{Winkler2023} established the global existence of generalized solution  
for arbitrary \( \alpha > 1 \). For the
sub-logistic source $f(n)$ satisfying \eqref{10.20-4} and 
\eqref{eq-1.9}, Xiang \cite{Xiang2018-2} proved that the sub-logistic source can 
prevent blow-up in 2D settings, which indicated that the logistic damping is not
the weakest damping that ensures boundedness for the model \eqref{eq:7.14-1}.

\subsection*{Keller-Segel-(Navier-)Stokes system with production of chemosignal}
Numerous cells in natural environments reside in viscous fluids and engage in 
significant mutual interactions with the surrounding fluid flow. Since fluid motion 
is governed by the incompressible (Navier-)Stokes equations, 
Tuval et al.\ \cite{Tuval-2005} proposed a coupled chemotaxis-fluid model:
\begin{equation}\label{10.20-2}
\begin{gathered}
n_t + u \cdot \nabla n = \Delta n -\xi\nabla\cdot(n \nabla v)  + f(n), \quad
 x \in \Omega\subset \mathbb{R}^N,\; t > 0, \\
v_t + u \cdot \nabla v = \Delta v - v + n, \quad
 x \in \Omega\subset \mathbb{R}^N,\; t > 0, \\
u_t +\kappa (u\cdot \nabla) u = \Delta u + \nabla P + n \nabla \phi +g, \quad
 x \in \Omega\subset \mathbb{R}^N,\; t > 0, \\
\nabla \cdot u = 0, \quad x \in \Omega\subset \mathbb{R}^N,\; t > 0.
\end{gathered}
\end{equation}
This model can be used to describe the dynamics of swimming aerobic bacteria, 
Bacillus subtilis, in water droplets. 
The coefficient $\kappa\geq 0$ is related to the intensity of nonlinear fluid 
convection. Specifically, $\kappa = 0$ means that the fluid flows slowly, and the
resulting system is called incompressible chemotaxis-Stokes system. 
If $ f(n)=0 $ and $ g=0 $ in \eqref{10.20-2}, Li and Xiao \cite{s1468} proved the 
global existence of the bounded classical solution provided that the initial mass
$\|n_0\|_{L^{1}(\Omega)} < \frac{\sqrt{6}}{2 C_{\mathrm{GN}}}$. 
Taking $ f(n)=-n^2$ and $g=0$ in \eqref{10.20-2}, Espejo and Suzuki \cite{Espejo} 
first gave the global existence of weak solution with $\kappa=0$ in the whole space
$\mathbb{R}^2$. Jin \cite{Jin} further derived the existence of strong solutions and 
analyzed their large-time behavior with $\kappa = 1$ in a bounded domain 
$\Omega\subset \mathbb{R}^2 $. In \cite{Zhang}, Zhang obtained the existence and 
uniqueness of weak solutions with $\kappa=1$ in $\mathbb{R}^2$ by employing the 
Fourier localization technique.

Tao and Winkler \cite{s1468-31} proposed a Keller-Segel-Navier-Stokes system of the 
form \eqref{10.20-2} with logistic damping $ f(n) = rn -\mu n^2 $ and 
$ g \neq 0 $ in a smooth bounded domain $\Omega \subset \mathbb{R}^3$. 
They proved that this system admits a global bounded classical solution provided that 
$\mu>23$ and $r\geq0$. Moreover, they also gave the large-time behavior of the 
solutions: when $r=0$, $\|n\|_{L^\infty(\Omega)}\rightarrow 0$, 
$\|c\|_{L^\infty(\Omega)}\rightarrow 0$ and $\|u\|_{L^\infty(\Omega)}\rightarrow 0$ 
as $t\rightarrow 0$. Later, Winkler \cite{Winkler2019}
showed that system \eqref{10.20-2} admits at least one global weak solution under the 
explicit hypothesis $\mu>\frac{\sqrt{r_+}}{4}$ ($r_+=\max\{0,r\}$). 
In addition, these solutions were proven to stabilize toward a spatially homogeneous 
steady state as time tends to infinity. For the 2D case, Tao and Winkler
\cite{s1468-30} found an interesting phenomenon that logistic damping 
$ f(n) = rn -\mu n^2 $ ($\mu>0$ and $r\geq0$) can prevent the blow-up of the solutions. 
Furthermore, Jin and Xiang \cite{Jin} refined the results in \cite{s1468-30} by giving 
how the upper bounds of solutions  qualitatively depend on $\chi$ and $\mu$. 
They also derived the large-time asymptotic behavior of the solutions under the 
assumption $r = 0$.

\subsection*{Keller-Segel-(Navier-)Stokes system with consumption of chemosignal} 
Another chemotaxis-(Navier-)Stokes system, different from the model \eqref{10.20-2}, 
with consumption of chemosignal was proposed and is formulated as 
\begin{equation}\label{eq:G-eq-1.11}
\begin{gathered}
n_t + u \cdot \nabla n = \Delta n -\nabla \cdot ( \chi(c) n \nabla c) + f(n), \quad
 x \in \Omega\subset \mathbb{R}^N,\; t > 0, \\
c_t + u \cdot \nabla c = \Delta c - h(c) n, \quad
 x \in \Omega\subset \mathbb{R}^N,\; t > 0, \\
u_t + \kappa ( u \cdot \nabla ) u = \Delta u + \nabla P + n \nabla \phi, \quad
 x \in \Omega\subset \mathbb{R}^N,\; t > 0, \\
\nabla \cdot u = 0, \quad x \in \Omega\subset \mathbb{R}^N,\; t > 0,
\end{gathered}
\end{equation}
where the concentration-dependent function $h(c)$ represents the oxygen consumption 
rate. We first recall some results related to the system \eqref{eq:G-eq-1.11} 
in the case $f(n)=0$. In \cite{G-8}, it was demonstrated that the
chemotaxis-Stokes system \eqref{eq:G-eq-1.11} admits globally defined classical 
solutions in the whole space $\mathbb{R}^3$, provided that the initial data 
$\|n_0, c_0, u_0\|_{H^3(\Omega)}$ is suitable small. Liu and Lord
\cite{s1468-17} continued to establish the global existence of the weak solution of 
the full chemotaxis-Navier-Stokes system only under
certain structural conditions governing the relationship between $\chi(c)$ and 
$h(c)$ in the whole space $\mathbb{R}^2$. For $ h(c) = c $, $ \chi(c) = \chi >0 $ and 
$ \kappa = 0 $, the system \eqref{eq:G-eq-1.11} reduces to the chemotaxis-fluid model,
which is proposed in \cite{Tuval-2005}. Due to the dissipative structure $-nc$, 
a new energy functional that incorporates the logarithmic entropy $\int_\Omega
n \ln n$ should be introduced. This constitutes a crucial step in deriving the 
global dynamical behavior of solutions, and it applies equally to
scenarios involving both the Stokes equation and the Navier-Stokes equation 
in \cite{W2020-28,W2020-29,W2020-30,W2020-31}.
In a series of seminal works, Winkler investigated the case $\chi(c) = 1 $ and 
$ h(c) = c $, establishing the following sharp results: for $N=2$, the full 
chemotaxis-Navier-Stokes system admits a unique global classical solution; 
for $N=3$, the simplified chemotaxis-Stokes system (with $\kappa=0 $)
possesses a global weak solution at least.

For the case $ f(n) \neq 0 $, $ \chi(c)\equiv\chi >0 $ and $ h(c) = c $, 
the model \eqref{eq:G-eq-1.11} has been studied by many authors. The results
concerning the case where the source term takes the common logistic form 
$ f(n) = n(1 - n) $ are established in \cite{ma-2,ma-14,ma-24}. In
particular, Baghaeia and Khelghatib \cite{ma-2} investigated the global existence and 
boundedness of solutions to the chemotaxis model under conditions 
$ \chi > 0$ and 
\[
 \|c_{0}\|_{L^{\infty}(\Omega)}
 \leq \frac{1}{\chi} \sqrt{\frac{d_1}{2(N+1)}} \Big[ \pi - 2\arctan \frac{d_1-1}{2}
\sqrt{\frac{2(N+1)}{d_1}}\Big].
\]
 Additionally, Xiang \cite{Xiang2018-2} demonstrated that a
sub-logistic source is capable of preventing blow-up in two-dimensional settings. 
Ma \cite{ma} proved the global existence and uniform-in-time boundedness of the 
solutions of the 2D chemotaxis-fluid model (with Neumann initial-boundary conditions), 
where the source term $f(n)$ has weaker damping source than
the standard logistic source.

\subsection*{Chemotaxis-fluid models with double chemical signals} 
In many natural biological processes, cells are influenced by more than one chemotactic
signal, each of which may act as an attractant or a repellent, giving rise to a 
variety of intricate patterns (see \cite{Painter}). In light of this,
Kozono,  Miura and Sugiyama \cite{kozono} proposed a chemotaxis-(Navier-)Stokes 
system with two different types of chemical attractants, formulated as 
\begin{equation}\label{eq:G-(1.8)}
\begin{gathered}
n_t + u \cdot \nabla n = \Delta n - \chi \nabla \cdot (n\nabla c) 
 + \xi \nabla \cdot (n\nabla v), \quad x \in \Omega\subset \mathbb{R}^N,\; t > 0, \\
c_t + u \cdot \nabla c = \Delta c - nc, \quad
 x \in \Omega\subset \mathbb{R}^N,\; t > 0, \\
v_t + u \cdot \nabla v= \Delta v - v + n, \quad
 x \in \Omega\subset \mathbb{R}^N,\; t > 0, \\
u_t + \kappa (u \cdot \nabla) u  = \Delta u + \nabla P + n\nabla \phi,  \quad
 x \in \Omega\subset \mathbb{R}^N,\; t > 0, \\
\nabla \cdot u = 0,  \quad x \in \Omega\subset \mathbb{R}^N,\; t > 0.
\end{gathered}
\end{equation}
This system presents a lot of the significant mathematical challenges due to its 
nature as a coupling of systems \eqref{10.20-2} and \eqref{eq:G-eq-1.11}. 
In \cite{kozono}, by using the implicit function theorem, Kozono,  Miura and Sugiyama 
constructed the unique global mild solution to the corresponding Cauchy problem in
$\mathbb{R}^N$ ($N \geq 2$) and its asymptotic behavior under some restrictive 
assumptions on the initial data. In the case where $\chi$ and $\xi$ are
positive constants, Ren \cite{G-35} proved the global existence of the 
chemotaxis-Navier-Stokes system \eqref{eq:G-(1.8)}. Recently, Xie and Xu
\cite{Global-1} established that the chemotaxis-Navier-Stokes system 
\eqref{eq:G-(1.8)} admits a global classical solution under
conditions $ \chi > 0$, $ \xi < 0$, 
$\|n_0\|_{L^1(\Omega)} < \min \{ \frac{1}{4| \xi | C_{\mathrm{GN}}}, 
\frac{2\pi}{ |\xi| } \}$ and
$\|c_0\|_{L^\infty(\Omega)} < \frac{1}{\chi}$. 

Inspired by the aforementioned works, we examine the chemotaxis-Stokes system 
with double chemical signals and sub-logistic source as
presented in \eqref{eq1.1}. We demonstrate that for any given real constant 
$ \chi \in \mathbb{R}$ and $ \xi\in \mathbb{R}$, system \eqref{eq1.1} admits a 
globally unique bounded classical solution. 
This result indicates that the sub-logistic source can prevent blow-up arising 
from the chemotaxis-Stokes system. Below, we state our main result.

\begin{theorem}\label{Th1.1}
Let $\Omega \subset \mathbb{R}^2$ be a bounded smooth domain, 
$\chi, \xi \in \mathbb{R}$, and let $(n_0, c_0, v_0, u_0)$ satisfy \eqref{eq1.3}. 
Suppose further that the potential function $\phi$ satisfies \eqref{eq1.2}, and
the source function $f\in W_{loc}^{1, \infty}(\mathbb{R}^+)$ satisfies 
$f(s)\geq \zeta s$ for small $s\geq0$, where $\zeta\in \mathbb{R}$, and the growth 
condition
\begin{equation}\label{10.20-4}
\liminf_{s \to \infty} \Big( - f(s) \frac{\ln s}{s^2} \Big) = \mu \in (0, \infty].
\end{equation}
If it holds that
\begin{equation}\label{eq-1.9}
\Big(2 + \chi^2 + \xi^2 + 4\|c_0\|_{L^\infty(\Omega)}^2
+ 4 C_{\mathrm{p}}^2\|c_0\|_{L^\infty(\Omega)}^2 \|\nabla \phi\|_{L^\infty(\Omega)}^2 
-\mu\Big)^+ < \frac{1}{2 C_{\mathrm{GN}}^4 M},
\end{equation}
where $C_{\mathrm{GN}} > 0$ denotes the Gagliardo-Nirenberg constant and 
$C_{\mathrm{p}}$ the Poincar\'e constant for $\Omega$, and 
\begin{equation}\label{eq-1.10}
M :=\|n_0\|_{L^1(\Omega)}+|\Omega| \inf_{\eta>0} 
\frac{\sup \{ f(s) + \eta s : s > 0 \}}{\eta},
\end{equation}
then the initial-boundary problem \eqref{eq1.1}-\eqref{eq1.3} admits a unique 
global-in-time solution $(n, c, v, u, P)$ satisfying
\begin{equation}\label{eq-1.11}
n, c, v, u \in C^0(\bar{\Omega} \times [0, \infty)) \cap C^{2, 1}(\bar{\Omega} \times (0, \infty)), P \in C^{1, 0}(\bar{\Omega} \times (0, \infty)).
\end{equation}
Moreover, there exists a constant $C > 0$ independent of $t$ such that
\begin{equation}\label{eq-1.12}
\| n(\cdot, t) \|_{L^\infty(\Omega)} + \| c(\cdot, t) \|_{W^{1,\infty}(\Omega)} 
+ \|v(\cdot, t)\|_{W^{1,\infty}(\Omega)} + \| u(\cdot, t)
\|_{L^\infty(\Omega)} \leq C  \quad \text{for all } t > 0.
\end{equation}
\end{theorem}

\begin{remark}\label{rem1.1} \rm
For sub-logistic sources like \( f(n) = n \bigl( a- \frac{bn}{\ln^\gamma(n+1)}\bigr)\) 
with \( a \in \mathbb{R}, b > 0, \gamma \in (0,1) \) or 
\( f(n) = n\bigl(a-\frac{bn}{\ln(\ln(n+e))}\bigr)\) with \( a \in \mathbb{R}, b>0\), 
one can easily compute from \eqref{10.20-4} that \( \mu= +\infty \) and so 
\eqref{eq-1.9} holds trivially.
\end{remark}

%\begin{remark}\label{re1.2} \rm
To the best of our knowledge, Theorem \ref{Th1.1} provides the first rigorous 
mathematical result concerning the double-chemical-signal chemotaxis-Stokes system 
\eqref{eq1.1} with a sub-logistic source. This theorem also extends the existing results, 
for instance, those in \cite{ma} where $v\equiv0$ in \eqref{eq1.1}, by providing an 
explicit condition \eqref{eq-1.9} that quantifies the relationship between 
$\mu$ (characterizing the sub-logistic source) and the physical coefficients (e.g., 
$\chi, \xi$).

%\begin{remark}\label{re1.3}
In contrast to existing results for chemotaxis-Navier-Stokes systems with nonlinear 
diffusion (see \cite{Luo}), where $\Delta n$ is replaced by $\nabla\cdot (D(n)\nabla n)$ 
and $f(n)=0$ (cf. \eqref{eq1.1}), our result in Theorem \ref{Th1.1} does not rely on 
the signs of $\chi$ and $\xi$ (under the condition \eqref{eq-1.9}). Furthermore, 
we establish not only the global existence but also the uniform-in-time boundedness 
of the classical solution, which constitutes an improvement over the aforementioned
 works.

This article is organized as follows. 
In Section 2, we recall some well-known lemmas, facts and inequalities that will be 
frequently used in the subsequent proofs. 
In Section 3, we first derive a series of a priori estimates for the local solution
 \((n,c,v,u)\) (whose existence is guaranteed by Lemma \ref{lm2.8}), and finally 
establish Theorem \ref{Th1.1} by applying the continuity method.

\section{Preliminaries}

In this section, we recall some preliminary lemmas that will be used frequently 
hereafter.
Firstly, we recall the well-known Gagliardo-Nirenberg inequality 
(see \cite{Nirenberg-1959}).

\begin{lemma}\label{lem2.1}
Let $ \Omega \subset \mathbb{R}^2 $ be a bounded smooth domain and let $ p \geq 1 $, 
$ r > 0 $.
\begin{itemize}
\item[(i)] For any $ q \in (0, p) $, there exists a
positive constant $ C_{\mathrm{GN}} = C(p, r, \Omega) $ such that
\begin{equation}\label{9.12-1}
\|\varphi\|_{L^p(\Omega)} \leq C_{\mathrm{GN}} ( \|\nabla \varphi\|_{L^2(\Omega)}^\theta \|\varphi\|_{L^q(\Omega)}^{1 - \theta} + \|\varphi\|_{L^r(\Omega)} )
\quad \text{for all } \varphi \in H^1(\Omega) \cap L^q(\Omega),
\end{equation}
where
\(\theta = 1 - \frac{q}{p} \in (0, 1)\).

\item[(ii)] For some $q, l, s \geq 1$, $j, m \in \mathbb{N}_0$ and 
$\theta \in [ \frac{j}{m},1]$ satisfying
\[
\frac{1}{p}=\frac{j}{2}+\Big( \frac{1}{l}-\frac{m}{2} \Big)\theta 
+ \frac{1-\theta}{q}.
\] 
Then there is a positive constant $C$ such that
\begin{equation}\label{l2.7.10.1}
\| D^j \varphi \|_{L^p(\Omega)} 
\leq C \big( \| D^m \varphi \|_{L^l(\Omega)}^\theta \| \varphi 
\|_{L^q(\Omega)}^{1-\theta} + \| \varphi \|_{L^s(\Omega)} \big)
\end{equation}
holds for all $\varphi \in W^{m,l}(\Omega) \cap L^s(\Omega)$.
\end{itemize}
\end{lemma}

To apply the integrability condition $|n \ln n|\in L^1(\Omega)$ for deriving further 
estimates (see the proof of Lemma \ref{lem3.5}), we need the following
logarithmic version of the generalized Gagliardo-Nirenberg inequality 
(see \cite{Winkler2014}).

\begin{lemma}\label{lem2.3}
Let $ \Omega \subset \mathbb{R}^2 $ be a bounded domain with smooth boundary, 
and let $ q >1 $ with $ r \in (0, q) $. Then, for each $
\varepsilon > 0 $ there exists a constant $ C_\varepsilon > 0 $ such that
\begin{equation}\label{9.12-2}
\|\varphi\|_{L^q(\Omega)}^{q} 
\leq \varepsilon \|\nabla \varphi\|_{L^2(\Omega)}^{q - r} \left\| |\varphi| 
\ln |\varphi| \right\|_{L^r(\Omega)}^{r} +
C_\varepsilon \|\varphi\|_{L^r(\Omega)}^{q} + C_\varepsilon
\end{equation}
holds for all $\varphi \in H^{1}(\Omega) \cap L^r(\Omega)$.
\end{lemma}

Next, we present two inequalities that are special cases of the trace theorem 
(see \cite{ma-13}).

\begin{lemma} \label{lem2.2}
Let $\Omega \subset \mathbb{R}^2$  be a bounded smooth domain and let 
$\varphi \in C^2(\bar{\Omega})$ satisfy the Neumann
boundary condition $\frac{\partial \varphi}{\partial \nu} = 0$
on $\partial \Omega$.
\begin{itemize}
\item[(i)] Then
\begin{equation}\label{9.12-3}
\frac{\partial}{\partial\nu }|\nabla \varphi|^2 
\leq \mathfrak{R} |\nabla \varphi|^2,
\end{equation}
where $\mathfrak{R}$ is an upper bound on the curvature of 
$\partial \Omega$.

\item[(ii)] Moreover, for any $\varepsilon > 0$ there is $C_\varepsilon> 0$ such that
\begin{equation}\label{9.12-4}
\|\nabla \varphi\|_{L^2(\partial \Omega)} \leq \varepsilon\|\Delta \varphi\|_{L^2(\Omega)}+C_\varepsilon \|\varphi\|_{L^1(\Omega)}.
\end{equation}
\end{itemize}
\end{lemma}

To derive some preliminary time-independent estimates, we recall an auxiliary lemma 
on boundedness of solutions to a linear differential inequalities,
which was established in \cite{ma-19}.

\begin{lemma}\label{lem2.4}
Let $T>0$, $\tau\in(0,T)$, $a > 0$ and $b>0$. Suppose that $y:[0,T)\to[0,\infty)$ 
is absolutely continuous and satisfies
\[
y'(t)+a y(t)\leq h(t)\quad\text{for a.e. }t\in(0,T).
\]
If the nonnegative function $h\in L^1_{loc}([0,T))$ satisfies
\[
\int_{t}^{t + \tau}h(s)ds\leq b\quad\text{for all }t\in[0,T-\tau),
\]
the following uniform bound holds for all $t\in(0,T)$
\[
y(t)\leq\max\big\{y(0)+b,\frac{b}{a\tau}+2b\big\}.
\]
\end{lemma}

In view of the central importance of semigroups in the following analysis, we recall 
some key estimates for the Neumann heat semigroup
$(e^{t\Delta})_{t\geq0}$ in the following lemma, the proof of which can be found 
in \cite{Cao2015, Horstmann2005,G-10}.

\begin{lemma}\label{lem2.6}
Let $(e^{t\Delta})_{t > 0}$ be the Neumann heat semigroup on a bounded smooth domain 
\(\Omega\subset \mathbb{R}^2\), and let $\lambda_1 > 0$ denote the first nonzero 
eigenvalue of $-\Delta$ in $\Omega$ under Neumann boundary condition. 
Then, there exist two constants $N_1>0$, $N_2>0$, and $N_3 > 0$, depending only on 
$\Omega$, such that:
\begin{itemize}
\item[(i)] For each $1 \leq q \leq p \leq \infty$, the estimate
\begin{equation}\label{9.12-6}
\|\nabla e^{t\Delta}  w\|_{L^p(\Omega)} 
\leq N_1 \Big(1 + t^{-\frac{1}{2}-\left(\frac{1}{q} - \frac{1}{p}\right)}\Big) 
e^{-\lambda_1 t} \|w\|_{L^q(\Omega)}
\end{equation}
holds for all $t > 0$ and for any $w \in L^{q}(\Omega)$.

\item[(ii)] For each \(2 \leq p < \infty\), the estimate
\begin{equation}\label{1.18-6}
\|\nabla e^{t\Delta} w\|_{L^p(\Omega)} 
\leq N_2 e^{-\lambda_1 t} \|\nabla w\|_{L^p(\Omega)}
\end{equation}
holds for all $t > 0$ and for any \(w \in W^{1,p}(\Omega)\).

\item[(iii)] For each $1 < q \leq p \leq \infty$, the estimate
\begin{equation}\label{9.12-6-2}
\|e^{t\Delta}\nabla \cdot w\|_{L^p(\Omega)} 
\leq N_3 \Big(1 + t^{-\frac{1}{2} - \left(\frac{1}{q} 
- \frac{1}{p}\right)}\Big) e^{-\lambda_1 t} \|w\|_{L^q(\Omega)}
\end{equation}
holds for all $t > 0$ and for any $w \in (W^{1,p}(\Omega))^n$.
\end{itemize}
\end{lemma}

Similarly, we need to introduce a semigroup to deal with the velocity $u$ 
in our system \eqref{eq1.1}. To this end,  we first recall the Helmholtz projection 
operator \(\mathcal{P}\), which maps the space $L^2(\Omega)$ onto its closed subspace
\[
L_\sigma^2(\Omega) := \{\varphi \in L^2(\Omega) \mid \nabla \cdot \varphi 
= 0 \quad \text{in } \mathcal{D}'(\Omega)\}.
\]
The Stokes operator $\mathcal{A} := -\mathcal{P}\Delta$ with domain 
$D(\mathcal{A}) = W^{2,p}(\Omega) \cap W_0^{1,p}(\Omega) \cap L_\sigma^2(\Omega)$
is sectorial in $L_\sigma^2(\Omega)$. Consequently, it generates the analytic Stokes 
contraction semigroup $(e^{-t\mathcal{A}})_{t\geq0}$ and admits
densely defined fractional powers $\mathcal{A}^\alpha$ for any 
$\alpha \in (0,1)$ (see \cite{Giga-1991, Sohr-2001}). 
Among the standard embedding and regularity estimates for 
\(\mathcal{A}\), we highlight the following lemma, whose proof can be derived by 
a minor modification of the three-dimensional counterpart presented in 
\cite[Lemma 3.1--3.3]{Winkler-2015}.

\begin{lemma}\label{lem2.7} 
There exists a constant $\lambda > 0$ such that for all $\eta \geq 0$ and $t\geq 0$, 
the estimate
\begin{equation}\label{eq-2.9}
    \|\mathcal{A}^\eta e^{-t\mathcal{A}} \mathcal\varphi\|_{L^p(\Omega)} \leq C(\eta) t^{-\eta} e^{-\lambda t}  \|\varphi\|_{L^p(\Omega)}
\end{equation}
 holds for all $\varphi \in L_\sigma^p(\Omega)$.
\end{lemma}

Finally, we state the local solvability and extendibility criterion for the 
chemotaxis-growth system \eqref{eq1.1}. These results can be rigorously established 
via an appropriate fixed-point argument combined with standard parabolic regularity 
theory (see, e.g., \cite{Winkler2007,Winkler2010}). For the sake of brevity, 
we omit the detailed proof.

\begin{lemma}\label{lm2.8}
Let $\chi, \xi\in \mathbb{R}$ be any given constants, and let
$\Omega \subset \mathbb{R}^2$ be a bounded domain with smooth boundary. 
Suppose that the initial data $(n_0, c_0, v_0 , u_0)$ satisfy
\eqref{eq1.3}, the potential $\phi$ satisfies \eqref{eq1.2}, and the source function 
$f \in W_{loc}^{1,\infty}(\mathbb{R}^+)$. There exists
a maximal existence time $T_{\rm max} \in (0, \infty]$ and a unique classical solution
$(n, c, v, u, P)$ to system \eqref{eq1.1} such that
\begin{gather*}
n \in C^0(\bar{\Omega} \times [0, T_{\rm max})) \cap C^{2, 1}(\bar{\Omega}
 \times (0, T_{\rm max})), \\
c \in \cap_{q > 2} C^0([0, T_{\rm max}); W^{1, q}(\Omega)) \cap C^{2, 1}(\bar{\Omega}
 \times (0, T_{\rm max})), \\
v \in \cap_{q > 2} C^0([0, T_{\rm max}); W^{1, q}(\Omega)) \cap C^{2, 1}(\bar{\Omega}
 \times (0, T_{\rm max})), \\
u \in \cap_{\alpha \in (\frac{1}{2}, 1)} C^0([0, T_{\rm max}); D(A^\alpha)) 
\cap C^{2, 1}(\bar{\Omega} \times (0, T_{\rm max}); \mathbb{R}^2),
\end{gather*}
and $n > 0$, $c > 0$ and $v \geq 0$ in $\bar{\Omega} \times (0, T_{\rm max})$. 
Moreover, for all $\alpha \in (\frac{1}{2}, 1)$, the following extendibility alternative 
holds: either $T_{\rm max} = +\infty$ or
\[
\limsup_{t \to T_{\rm max}^-} ( \|n(\cdot, t)\|_{L^\infty(\Omega)} 
+ \| c(\cdot, t) \|_{W^{1,\infty}(\Omega)} + \|v(\cdot, t)\|_{W^{1,\infty}(\Omega)} +
\|A^\alpha u(\cdot, t)\|_{L^2(\Omega)} ) = \infty.
\]
\end{lemma}

\section{Proof of Theorem \ref{Th1.1}}

In this section, we focus on the existence and boundedness of global  solutions 
to system \eqref{eq1.1}. To this end, we shall derive a series of a priori estimates.
We first establish the \(L^1\)-estimates for \(n\) and \(v\). These estimates are 
particularly crucial due to the strong nonlinearity of the equation governing \(n\); 
the low integrability of \(L^1\) introduces considerable challenges to the subsequent 
analysis.

\begin{lemma}\label{lem3.1}
Under the conditions in Theorem \ref{Th1.1}, we have
\begin{gather}\label{10.20-3}
\|n(\cdot,t)\|_{L^1(\Omega)} \leq M \quad \text{for all } t \in (0, T_{\rm max}), \\
\label{eq-3.2}
\|v(\cdot,t)\|_{L^1(\Omega)} \leq \|v_0\|_{L^1(\Omega)}+M  
\quad \text{for all } t \in (0, T_{\rm max}),
\end{gather}
where the constant $M$ is given in \eqref{eq-1.10}.
\end{lemma}

\begin{proof} 
To prove \eqref{10.20-3}, we first analyze the source term $f(n)$. By the definition 
$\mu$ in \eqref{10.20-4}, there exists a positive
constant \( \hat{s} \geq 1 \) such that
\begin{equation}\label{eq-3.3}
f(s) \leq - \hat{\mu} \frac{s^2}{\ln s} \quad \text{for all } s \geq \hat{s},
\end{equation}
where $\hat{\mu}  = \frac{\mu}{2}$ if $0 < \mu < \infty$ and $\hat{\mu}$  
can be chosen as any constant greater than $\chi$ if $\mu=\infty$.
For any $\eta>0$, it holds that
\begin{equation}\label{eq3.4}
f(s) + \eta s \leq -\hat{\mu} \frac{s^2}{\ln s} + \eta s \quad 
\text{for all } s \geq \hat{s}.
\end{equation}
This, together with the fact  $f$ is bounded on any finite interval, 
yields for any $\eta>0$,
\[
M_\eta :=  \sup\left\{ f(s) + \eta s : s > 0 \right\} < \infty.
\]

Integrating \eqref{eq1.1}$_1$ over $\Omega$ and using the homogeneous Neumann 
boundary conditions, we obtain that for any $\eta > 0$,
\begin{equation}\label{eq-3.5}
\frac{d}{dt} \int_{\Omega} n = \int_{\Omega} f(n) \leq -\eta \int_{\Omega} n + M_\eta |\Omega|,
\end{equation}
which, together with \eqref{eq-1.10}, implies that
\[
\int_{\Omega} n \leq \int_{\Omega} n_0 + \frac{M_\eta}{\eta} |\Omega| \leq M.
\]
This is the desired estimate \eqref{10.20-3}.

Since $\nabla\cdot u = 0$, by integrating \eqref{eq1.1}$_3$ and using \eqref{10.20-3}, 
we have
\[
\frac{d}{dt}\int_{\Omega} v + \int_{\Omega} v 
= \int_{\Omega} n \leq M \quad \text{for all } t \in (0, T_{\rm max}).
\]
Solving the above differential inequality directly yields
\[
\int_{\Omega} v \leq \|v_0\|_{L^1(\Omega)}+M\quad \text{for all } 
t \in (0, T_{\rm max}),
\]
which is exactly estimate \eqref{eq-3.2}. 
This completes the proof.
\end{proof}

Next, we apply the parabolic smoothing property of the momentum equation 
\eqref{eq1.1}$_4$ to establish $L^p$-estimate of $u$.

\begin{lemma}\label{lm3.2}
For any $ r \in [1, 2) $, there exists a positive constant $K_1$ such that
\begin{align}\label{10.20-1}
\| Du (\cdot, t) \|_{L^r(\Omega)} \leq K_1
\quad \text{for all } t \in (0, T_{\rm max}).
\end{align}
Furthermore, for each $p\in [1,+\infty)$, there exists a positive constant $K_2$ 
such that
\begin{equation}\label{eq3.6}
\|u(\cdot, t)\|_{L^p(\Omega)} \leq K_2 \quad \text{for all } t \in (0, T_{\rm max}).
\end{equation}
\end{lemma}

\begin{proof} 
Based on the conditions \eqref{eq1.2} and the $L^1$-estimate \eqref{10.20-3} for $n$, 
estimate \eqref{10.20-1} is a direct consequence of \cite[Lemma 2.5]{Wang-Xiang}. 
Furthermore, by the Sobolev embedding theorem 
$W^{1,r}(\Omega)\hookrightarrow L^p(\Omega)$ (valid for $r\in[1,2)$ and any 
$p\in[1,+\infty)$ in two dimensional domains), one easily obtains estimate \eqref{eq3.6}. 
This completes the proof.
\end{proof}

By the non-negativity of $n$ and $c$, we derive the next $L^p$-estimate for $c$.

\begin{lemma}\label{lem3.4}
Let $1 \leq p <\infty $. Then
\begin{align}\label{eq-3.8}
\|c(\cdot,t)\|_{L^p(\Omega)} \leq \|c_0\|_{L^p(\Omega)} \quad 
\text{for all } t \in (0, T_{\rm max}).
\end{align}
In particular, it also holds that
\begin{align}\label{(9.12.3.15)}
\|c(\cdot,t)\|_{L^\infty(\Omega)} \leq \|c_0\|_{L^\infty(\Omega)} \quad 
\text{for all } t \in (0, T_{\rm max}).
\end{align}
\end{lemma}

\begin{proof} 
For any \(p \geq 1\), we multiply \eqref{eq1.1}$_2$ by \(c^{p - 1}\) and integrate 
the resulting equation over \(\Omega\). Using integration by parts and the homogeneous 
Neumann boundary condition, we obtain
\begin{align}\label{eq-3.10}
\frac{1}{p} \frac{d}{dt} \int_{\Omega} c^p &= -\int_{\Omega} c^{p - 1} u \cdot \nabla c + \int_{\Omega} c^{p-1} \Delta c-\int_{\Omega} n c^p
\notag\\
&= \frac{1}{p} \int_{\Omega} c^p \nabla \cdot u - (p - 1) \int_{\Omega} c^{p - 2} |\nabla c|^2 - \int_{\Omega} n c^p.
\end{align}
Recalling that $n > 0$, $c > 0$ and the incompressibility condition $\nabla \cdot u = 0$ in $\Omega \times (0, T_{\rm max})$, we deduce from \eqref{eq-3.10} that
\[
 \frac{d}{dt} \int_{\Omega} c^p \leq 0.
\]
Integrating the above inequality with respect to \(t\) immediately yields estimate 
\eqref{eq-3.8}.
Furthermore, letting $ p \rightarrow \infty$ in \eqref{eq-3.8}, this is the 
desired estimate \eqref{(9.12.3.15)}. The proof of Lemma \ref{lem3.4} is
complete. 
\end{proof}


\begin{lemma}\label{lem3.5}
There exists a constant $K_3>0$ such that for all $ t \in (0, T_{\rm max}) $
\begin{equation} \label{eq-3.11}
\|n \ln n (\cdot,t)\|_{L^1(\Omega)} + \|\nabla c(\cdot,t)\|_{L^2(\Omega)}
+ \|\nabla v(\cdot,t)\|_{L^2(\Omega)} + \| u(\cdot,t)\|_{L^2(\Omega)} \leq K_3.
\end{equation}
Moreover, there exists a constant $K_4$ such
\begin{equation} \label{3.18-1}
\int_{t}^{t+t_1}\int_{\Omega}\left(|\nabla n^{1/2}|^2+|\Delta c|^2+|\Delta v|^2
+|\nabla u|^2\right) \le K_4\quad \forall t\in(0,T_{\rm max}-t_1),
\end{equation}
where $t_1:=\min\left\{1,\frac{1}{6}T_{\rm max}\right\}$.
\end{lemma}

\begin{proof} 
We multiply \eqref{eq1.1}$_1$  by $ \ln n + 1 $ and integrate the resulting equation 
over $ \Omega $ to see that
\begin{equation}\label{(3.22)}
\begin{aligned}
&\frac{d}{dt} \int_{\Omega} n \ln n  + 4 \int_{\Omega} |\nabla n^{1/2}|^2
\\&= \int_{\Omega} u \cdot \nabla (n\ln n) 
 + \chi \int_{\Omega} \nabla n \nabla c-\xi \int_{\Omega} \nabla n \nabla v 
 + \int_{\Omega} (\ln n + 1)
f(n) \\
&= -\chi \int_{\Omega} n \Delta c +\xi \int_{\Omega} n \Delta v
  + \int_{\Omega} (\ln n + 1) f(n) \\
&\leq \frac{1}{4} \int_{\Omega} \left(|\Delta c|^2+|\Delta v|^2\right)
  + (\chi^2+\xi^2 ) \int_{\Omega} n^2 + \int_{\Omega} \left(\ln n + 1\right) f(n),
\end{aligned}
\end{equation}
where we used Young's inequality and the fact $\nabla \cdot u = 0$ in 
$\Omega \times (0, T_{\rm max})$. Following an analogous approach to the derivation 
of \eqref{(3.22)}, we multiply \eqref{eq1.1}$_2$ by $-2\Delta c$ and integrate the 
resulting equation by parts over $\Omega$. This, together with \eqref{(9.12.3.15)}, 
gives 
\begin{equation}\label{(3.23)}
\begin{aligned}
\frac{d}{dt} \int_\Omega |\nabla c|^2 + 2 \int_\Omega |\Delta c|^2
&= 2 \int_\Omega nc\Delta c + 2 \int_\Omega u \cdot \nabla c \Delta c \\
&\leq \frac{1}{4} \int_\Omega |\Delta c|^2 
 + 4\|c_0\|_{L^\infty(\Omega)}^2 \int_\Omega n^2 
 + 2 \int_\Omega u \cdot \nabla c \Delta c.
\end{aligned}
\end{equation}
To proceed further, we need to bound the last term on the right hand of \eqref{(3.23)}. 
Using the incompressibility condition $\nabla \cdot u=0$, the fact 
$u \cdot \nabla c =\nabla \cdot (uc)$, and the boundary condition, we obtain after
 using integration over $\Omega$ and Young's inequality that
\begin{equation}\label{(3.24)}
\begin{aligned}
2 \int_\Omega u \cdot \nabla c \Delta c
&= 2 \sum_{i,j} \int_\Omega \partial_j (u_j c)\partial_{ii} c
 = 2 \sum_{i,j} \int_\Omega \partial_i (u_j c)\partial_{ij} c \\
&= 2 \sum_{i,j} \int_\Omega c\partial_i u_j  \partial_{ij} c 
 + 2 \sum_{i,j} \int_\Omega u_j \partial_i c \partial_{ij} c  \\
&= 2 \int_\Omega c D^2 c : \nabla u -\int_\Omega \nabla\cdot u |\nabla c|^2 
 +\int_{\partial\Omega} |\nabla c|^2 u\cdot \nu \\
&= 2 \int_\Omega c D^2 c : \nabla u \leq 2\|c_0\|_{L^\infty(\Omega)}^2 
 \int_\Omega |\nabla u|^2 + \frac{1}{2} \int_\Omega |D^2 c|^2.
\end{aligned}
\end{equation}
Next, we establish a priori bound for $\int_\Omega |D^2 c|^2$ in terms of 
$\int_\Omega |\Delta c|^2$. Using integration by parts over
$\Omega$, \eqref{9.12-3}, \eqref{9.12-4} and \eqref{(9.12.3.15)}, we obtain, 
for all $\varepsilon>0$ and some $C_1$, $C_2>0$, that
\begin{equation}\label{(3.25)}
\begin{aligned}
\frac{1}{2} \int_\Omega |D^2 c|^2
&= \frac{1}{2} \int_\Omega |\Delta c|^2 
 + \frac{1}{4} \int_{\partial\Omega} \frac{\partial}{\partial \nu} |\nabla c|^2 \\
&\leq \frac{1}{2} \int_\Omega |\Delta c|^2 + C_1 \int_{\partial\Omega} |\nabla c|^2 \\
&\leq \frac{1}{2} \int_\Omega |\Delta c|^2 + C_1 \varepsilon \int_\Omega |\Delta c|^2 + C_\varepsilon \int_\Omega c^2\\
&\leq ( \frac{1}{2}+ C_1 \varepsilon ) \int_\Omega |\Delta c|^2 + C_2.
\end{aligned}
\end{equation}
Putting \eqref{(3.23)}, \eqref{(3.24)} and \eqref{(3.25)} together, and taking 
$\varepsilon$ satisfying $ C_1 \varepsilon = \frac{1}{4}$, we obtain
\begin{equation}\label{(3.26)}
\frac{d}{dt} \int_\Omega |\nabla c|^2 + \int_\Omega |\Delta c|^2
\leq 4\|c_0\|_{L^\infty(\Omega)}^2  \int_\Omega n^2 + 2 \|c_0\|_{L^\infty(\Omega)}^2  \int_\Omega |\nabla u|^2 + C_2.
\end{equation}

We multiply \eqref{eq1.1}$_3$ by $2\Delta v$ in $L^2$ and using integration
 by parts and Young's inequality to deduce that
\begin{equation}\label{(3.28)}
\begin{aligned}
\frac{d}{dt} \int_{\Omega} |\nabla v|^2 
+ 2 \int_{\Omega} |\nabla v|^2 + 2 \int_{\Omega} |\Delta v|^2 
&= - 2 \int_{\Omega} n \Delta v +
2\int_{\Omega} \Delta v (u \cdot \nabla v) \\
&\leq \int_{\Omega} |\Delta v|^2 + \int_{\Omega} n^2 
+ 2\int_{\Omega} \Delta v (u \cdot \nabla v) .
\end{aligned}
\end{equation}
Using \eqref{l2.7.10.1}, \eqref{eq-3.2}, \eqref{10.20-1} and Young's inequality, we have
\begin{equation}\label{eq-3.19}
\begin{aligned}
2\int_{\Omega} \Delta v (u \cdot \nabla v) 
&= -2\sum_{i,j=1}^{2} \int_{\Omega} \partial_i u^j \partial_i v \partial_j v+\int_{\Omega}|\nabla
v|^2\nabla\cdot u\\
&\leq C_3 \|\nabla u\|_{L^{\frac{5}{4}}(\Omega)} \|\nabla v\|_{L^{10}(\Omega)}^2 \\
&\leq C_4\left(\|v\|_{L^1(\Omega)}^{\frac{2}{15}} \|\Delta v\|_{L^2(\Omega)}^{\frac{28}{15}}+\| v\|^2_{L^1(\Omega)}\right) \\
&\leq \frac{1}{2}\|\Delta v\|_{L^2(\Omega)}^2+C_5.
\end{aligned}
\end{equation}
Then, substituting the above inequality into \eqref{(3.28)}, yields
\begin{equation}\label{eq-3.20}
\frac{d}{dt} \int_{\Omega} |\nabla v|^2 + 2 \int_{\Omega} |\nabla v|^2 
+ \frac{1}{2} \int_{\Omega} |\Delta v|^2 \leq \int_{\Omega} n^2+ C_6.
\end{equation}
Next, we multiply \eqref{eq1.1}$_3$ by $2v$ and use integration by parts and 
Young's inequality to deduce the differential inequality
\begin{equation}\label{eq-3.21}
\frac{d}{dt} \int_{\Omega} v^2 + 2\int_{\Omega} |\nabla v|^2
 + 2\int_{\Omega} v^2 = 2\int_{\Omega} nv \leq \int_{\Omega} n^2 + \int_{\Omega} v^2.
\end{equation}
Putting \eqref{eq-3.21} and \eqref{eq-3.20} together leads to
\begin{equation}\label{eq-3.22}
\frac{d}{dt} \int_{\Omega} ( v^2 + |\nabla v|^2 ) 
+ \int_{\Omega}\Big(\frac{1}{2} |\Delta v|^2+4|\nabla v|^2+ v^2\Big) 
\leq 2\int_{\Omega} n^2+ C_6.
\end{equation}

Again, we multiply \eqref{eq1.1}$_4$ by $u$ in $L^2$ and use integration by parts 
and Young's inequality to deduce the differential inequality
\begin{equation}\label{eq-3.23}
\begin{aligned}
\frac{d}{dt} \int_{\Omega} u^2 + 2\int_{\Omega} |\nabla u|^2
&= 2\int_{\Omega} n \nabla \phi \cdot u\leq 2\|\nabla\phi\|_{L^\infty}\|u\|_{L^2}\|n\|_{L^2}
\\
& \leq \int_\Omega |\nabla u|^2 + C_{\mathrm{p}}^2\|\nabla\phi\|^2_{L^{\infty}(\Omega)}
\int_\Omega n^2. 
\end{aligned}
\end{equation}
where we used Poincar\'e inequality to estimate $\|u\|_{L^2}$ as
$$
\|\omega\|_{L^2}\leq C_{\mathrm{p}}\|\nabla \omega\|_{L^2}\text{for } 
\omega\in W^{1,2}_0(\Omega),
$$
where $C_{\mathrm{p}}$ is the Poincar\'e constant.

Then, we multiply \eqref{eq-3.23} by $4\|c_0\|^2_{L^\infty(\Omega)}$ and sum 
the resulting inequality with \eqref{(3.22)}, \eqref{(3.26)} and \eqref{eq-3.22}. 
This combination yields the key differential inequality
\begin{equation}\label{eq-3.24}
\begin{aligned}
&\frac{d}{dt} \int_{\Omega} \left( n \ln n - n + 1 + |\nabla c|^2 
+ |\nabla v|^2 + 4\|c_0\|_{L^\infty(\Omega)}^2 u^2 +  v^2\right) \\
& +\int_{\Omega}\Big(\frac{1}{4} |\Delta v|^2+4|\nabla v|^2+ v^2\Big) 
+ 4 \int_{\Omega}| \nabla n^{1/2}|^2
+ \frac{3}{4} \int_{\Omega} |\Delta c|^2  
+2 \|c_0\|_{L^\infty(\Omega)}^2 \int_{\Omega} |\nabla u|^2 \\
&\leq \left( 2 + \chi^2 + \xi^2+4\|c_0\|_{L^\infty(\Omega)}^2 
 + 4 C_{\mathrm{p}}^2\|c_0\|_{L^\infty(\Omega)}^2 
 \|\nabla \phi\|_{L^\infty(\Omega)}^2\right) 
 \int_{\Omega} n^2 + \int_{\Omega}  f(n)\ln n + C_7,
\end{aligned}
\end{equation}
where \eqref{eq-3.5} is also used.
By  \eqref{10.20-4}, for any $\varepsilon \in (0, \mu)$ there exists a constant 
$\hat{s}_0\geq 1$ such that
\begin{equation} \label{eq-3.25}
    f(s) \leq -(\mu - \varepsilon) \frac{s^2}{\ln s} \ \text{for all } s > \hat{s}_0.
\end{equation}
Then we have
\begin{align*}
&(2 + \chi^2 + \xi^2+  4\|c_0\|_{L^\infty(\Omega)}^2
 +4 C_{\mathrm{p}}^2\|c_0\|_{L^\infty(\Omega)}^2 \|\nabla \phi\|_{L^\infty(\Omega)}^2 ) s^2 +  f(s) \ln s \\
&\leq ( 2 + \chi^2 + \xi^2+ 4\|c_0\|_{L^\infty(\Omega)}^2
 +4 C_{\mathrm{p}}^2\|c_0\|_{L^\infty(\Omega)}^2 \|\nabla \phi\|_{L^\infty(\Omega)}^2) s^2 - (\mu - \varepsilon)s^2\\
&\leq \left( 2 + \chi^2 + \xi^2 + 4\|c_0\|_{L^\infty(\Omega)}^2
 + 4 C_{\mathrm{p}}^2\|c_0\|_{L^\infty(\Omega)}^2 \|\nabla \phi\|_{L^\infty(\Omega)}^2 - \mu\right)^+ s^2 + \varepsilon s^2 \quad
\text{for all } s \geq \hat{s}_0.
\end{align*}
This, together with \eqref{eq-3.24}, leads to
\begin{equation}\label{eq-3.26}
\begin{aligned}
&\frac{d}{dt} \int_{\Omega} \left( n \ln n - n + 1 + |\nabla c|^2 + |\nabla v|^2 + 4\|c_0\|_{L^\infty(\Omega)}^2 u^2 + v^2 \right)\\
& +\int_{\Omega}\Big(\frac{1}{4} |\Delta v|^2+4|\nabla v|^2+ v^2\Big)
 + 4 \int_{\Omega}| \nabla n^{1/2}|^2 + \frac{3}{4} \int_{\Omega} |\Delta c|^2
 +2 \|c_0\|_{L^\infty(\Omega)}^2 \int_{\Omega} |\nabla u|^2\\
&\leq \int_{\{n \leq \hat{s}_0\}} \left[ (2 + \chi^2 + \xi^2
 + 4\|c_0\|_{L^\infty(\Omega)}^2+ 4 C_{\mathrm{p}}^2\|c_0\|_{L^\infty(\Omega)}^2
 \|\nabla \phi\|_{L^\infty(\Omega)}^2 )n^2 +  f(n) \ln n \right] \\
&\quad + \int_{\{n >\hat{s}_0\}} \left[ \left(2 + \chi^2 + \xi^2
 + 4\|c_0\|_{L^\infty(\Omega)}^2+ 4 C_{\mathrm{p}}^2\|c_0\|_{L^\infty(\Omega)}^2
 \|\nabla \phi\|_{L^\infty(\Omega)}^2 \right)n^2 +  f(n)
\ln n \right]+C_7 \\
&\leq \Big[\left(2 + \chi^2 + \xi^2 + 4\|c_0\|_{L^\infty(\Omega)}^2
 + 4 C_{\mathrm{p}}^2\|c_0\|_{L^\infty(\Omega)}^2 \|\nabla \phi\|_{L^\infty(\Omega)}^2
 - \mu\right)^+ + \varepsilon\Big] \int_{\Omega} n^2 + C_8,
\end{aligned}
\end{equation}
where we have used the assumption $f(s)\geq \zeta s$ for small $s\geq0$, where
$\zeta\in \mathbb{R}$, to deduce that
$$
\int_{\{n \leq \hat{s}_0\}} f(n)\ln n\leq C_9,
$$
for some positive constant $C_9$.

Using \eqref{9.12-1}, \eqref{10.20-3}, \eqref{eq-3.8} and Young's inequality, we obtain
\begin{equation}\label{9.14-1}
\|\nabla c(t)\|_{L^2(\Omega)}^2
\leq C_{10} \|\Delta c(t)\|_{L^2(\Omega)} \|c(t)\|_{L^2(\Omega)}
+ C_{10} \|c(t)\|^2_{L^2(\Omega)} 
\leq \frac{1}{2} \|\Delta c(t)\|_{L^2(\Omega)}^2 + C_{11},
\end{equation}
and
\begin{equation}\label{eq-3.28}
\begin{aligned}
\int_{\Omega} n^2 = \| n^{1/2}\|_{L^4(\Omega)}^4
&\leq C_{\mathrm{GN}}^4 ( \| \nabla n^{1/2} \|_{L^2(\Omega)}^{1/2}
\| n^{1/2} \|_{L^2(\Omega)}^{1/2} + \| n^{1/2}\|_{L^2(\Omega)} )^4 \\
&\leq 8 C_{\mathrm{GN}}^4 ( M \| \nabla n^{1/2} \|_{L^2(\Omega)}^2 + M^2),
\end{aligned}
\end{equation}
because $ (a + b)^4 \leq 8(a^4 + b^4) $ for all $ a, b \geq 0$.
Next, noticing that
\[
n\ln n \leq \varepsilon n^2 + L_\varepsilon,\quad 
L_\varepsilon = \sup\left\{ s \ln s - \varepsilon s^2 : s > 0 \right\}<\infty,
\]
we conclude from \eqref{eq-3.28} that for any $\varepsilon>0$
\begin{equation}\label{3.40-1}
\int_{\Omega} n\ln n \leq 8\varepsilon M C_{\mathrm{GN}}^4
\int_{\Omega} |\nabla n^{\frac12}|^2 + C_{12}.
\end{equation}
Then, summing  \eqref{eq-3.26}, \eqref{9.14-1}, \eqref{eq-3.28} and \eqref{3.40-1}, 
and using Poincar\'e inequality, we obtain
\begin{equation}\label{eq-3.30}
\begin{aligned}
&\frac{d}{dt} \int_\Omega \left( n \ln n - n + 1 + |\nabla c|^2 
 + |\nabla v|^2 + v^2+ 4\|c_0\|_{L^\infty(\Omega)}^2 u^2 \right)
 +\int_\Omega \left( n \ln n - n + 1\right)\\ 
&+\int_{\Omega} |\nabla c|^2 +\frac{1}{4}\int_{\Omega} |\Delta c|^2
 +\int_{\Omega}\Big(\frac{1}{4} |\Delta v|^2+4|\nabla v|^2+ v^2\Big) \\
&+\frac{\|c_0\|_{L^\infty(\Omega)}^2}{C^2_{\mathrm{p}}}\int_{\Omega} |u|^2
 +\|c_0\|_{L^\infty(\Omega)}^2\int_{\Omega} |\nabla u|^2  
 + 4 \Big( 1 - 2 C_{\mathrm{GN}}^4 M \Big[ \Big( 2 + \chi^2 + \xi^2 \\
&+ 4\|c_0\|_{L^\infty(\Omega)}^2+ 4 C_{\mathrm{p}}^2\|c_0\|_{L^\infty(\Omega)}^2
  \|\nabla \phi\|_{L^\infty(\Omega)}^2 - \mu \Big)^+ 
  + 2\varepsilon \Big] \Big) \int_{\Omega} | \nabla n^{1/2} |^2 \\
&\leq  C_{13}.
\end{aligned}
\end{equation}
Thanks to condition \eqref{eq-1.9}, we can fix a positive constant $\varepsilon$ 
in \eqref{eq-3.30} by setting
\[
\varepsilon = \frac{1}{4} \min \Big\{ \mu, \frac{1}{2 M C_{\mathrm{GN}}^4} 
- \Big(2 + \chi^2 + \xi^2 + 4\|c_0\|_{L^\infty(\Omega)}^2
 + 4 C_{\mathrm{p}}^2\|c_0\|_{L^\infty(\Omega)}^2 \|\nabla \phi\|_{L^\infty(\Omega)}^2 
 - \mu\Big)^+ \Big\}>0,
\]
where the positivity of $\varepsilon$ is guaranteed by \eqref{eq-1.9}. 
This, together with \eqref{eq-3.30}, yields that
\begin{align*}
&\frac{d}{dt} \int_{\Omega} \left( n \ln n - n + 1 + |\nabla c|^2 + |\nabla v|^2 + v^2+ 4\|c_0\|_{L^\infty(\Omega)}^2 u^2\right) \\
&+ \min\big\{ 1, \frac{1}{4 C^2_{\mathrm{p}}}\big\} 
\int_{\Omega} \Big(n \ln n - n + 1 + |\nabla c|^2 + |\nabla v|^2 + v^2
 + 4\|c_0\|_{L^\infty(\Omega)}^2 u^2\Big)
\leq C_{13}.
\end{align*}
Then, applying Lemma \ref{lem2.4} to above inequality leads to
\[
\int_\Omega \left(n \ln n - n + 1 + |\nabla c|^2 + |\nabla v|^2 + v^2
+ 4\|c_0\|_{L^\infty(\Omega)}^2 u^2\right)\leq C_{14}.
\]
This, together with the fact $-s\ln s\leq e^{-1}$ for $s>0$ and \eqref{10.20-3}, 
immediately leads to \eqref{eq-3.11}.

By  \eqref{eq-3.11}, integrating \eqref{eq-3.30} with respect to time leads to
\begin{equation} \label{3.18-2}
\int_{t}^{t+t_1}\int_{\Omega}\left(|\nabla n^{1/2}|^2
+|\Delta c|^2+|\Delta v|^2+|\nabla u|^2\right)
\le C_{15} \quad\text{for all } t\in(0,T_{\rm max}-t_1).
\end{equation}
This is the estimate given in \eqref{3.18-1}.
The proof is complete.
\end{proof}

Now, we are in a position to improve the integrability of $n$ by the above lemma.

\begin{lemma}\label{lem3.6}
There exist two positive constants $K_5$ and $K_6$ such that
\begin{gather}\label{eq-3.32}
\int_{\Omega} \left(n^2(\cdot, t)+| \nabla c(\cdot, t) |^4
+ | \nabla v(\cdot, t) |^4\right) \leq K_5 \quad\text{for all }
t \in (0, T_{\rm max}), \\
\label{1-18-2}
\int_t^{t+t_1}\int_{\Omega}\left(|\nabla n|^2+|\nabla |\nabla c|^2|^2
+|\nabla |\nabla v|^2|^2\right)\leq K_{6} \quad\text{for all } 
t \in (0, T_{\rm max}-t_1),
\end{gather}
where $t_1:=\min\{1,\frac{1}{6}T_{\rm max}\}$.
\end{lemma}

\begin{proof} 
We multiply \eqref{eq1.1}$_1$ by $2n$, then integrate by parts over $\Omega$ and 
apply Young's inequality to deduce that
\begin{equation}\label{eq-3.34}
\begin{aligned}
&\frac{\mathrm{d}}{\mathrm{d}t} \int_{\Omega} n^2 
+ 2\int_{\Omega} | \nabla n |^2
\\
&=2\chi \int_{\Omega} n \nabla n \cdot \nabla c - 2\xi \int_{\Omega} n \nabla n \cdot \nabla v + 2\int_{\Omega} n f(n)\\&\leq \int_{\Omega} |\nabla
n|^2 +\varepsilon^{-1/2}\int_{\Omega} n^3
+ C_1\varepsilon\int_{\Omega} |\nabla c|^6 
+ C_2\varepsilon\int_{\Omega} |\nabla v|^6+ 2\int_{\Omega} nf(n).
\end{aligned}
\end{equation}

After applying the gradient operator $\nabla$ to \eqref{eq1.1}$_3$ and taking the 
$L^2$-inner product of the resulting equation with $4|\nabla v|^2\nabla v$, we obtain
\begin{equation}\label{(3.45)}
\begin{aligned}
\frac{1}{4} \frac{d}{dt}  \int_{\Omega} | \nabla v |^4
&= \int_{\Omega} |\nabla v|^2 \nabla v \cdot \nabla \bigl( \Delta v - v + n - u \cdot \nabla v \bigr) \\
&= \frac{1}{2} \int_{\Omega} |\nabla v|^2 \Delta |\nabla v|^2
  - \int_{\Omega} |\nabla v|^2 |D^2 v|^2
  - \int_{\Omega} |\nabla v|^4 \\
&\quad- \int_{\Omega} n \nabla \cdot \bigl( |\nabla v|^2 \nabla v \bigr)
  + \int_{\Omega} (u \cdot \nabla v) \nabla \cdot \bigl( |\nabla v|^2 \nabla v \bigr) \\
&= -\frac{1}{2} \int_{\Omega} \bigl| \nabla |\nabla v|^2 \bigr|^2- \int_{\Omega} |\nabla v|^4 - \int_{\Omega} |\nabla v|^2 |D^2 v|^2 \\
&\quad + \frac{1}{2} \int_{\partial \Omega} |\nabla v|^2 \frac{\partial |\nabla v|^2}{\partial \nu}
  - \int_{\Omega} n |\nabla v|^2 \Delta v
  - \int_{\Omega} n \nabla v \cdot \nabla |\nabla v|^2\\
&\quad+ \int_{\Omega} (u \cdot \nabla v) |\nabla v|^2 \Delta v
  + \int_{\Omega} (u \cdot \nabla v) \nabla v \cdot \nabla |\nabla v|^2\\&:=-\frac{1}{2} \int_{\Omega} \bigl| \nabla |\nabla v|^2 \bigr|^2-
  \int_{\Omega} |\nabla v|^4 - \int_{\Omega} |\nabla v|^2 |D^2 v|^2 +\sum_{i=1}^{5}I_i,
\end{aligned}
\end{equation}
where the fact $ 2\nabla v \cdot \nabla \Delta v = \Delta | \nabla v |^2
- 2 | D^2 v |^2 $ is used.
Next, we proceed to estimate each term on the right hand of \eqref{(3.45)}. 
By \eqref{9.12-3}, \eqref{9.12-4}, and \eqref{eq-3.11}, we obtain
\begin{align*}
I_1 &\leq \frac{1}{2}\int_{\partial \Omega} |\nabla v|^4  \\ &\leq \frac{1}{16}\int_{\Omega} |\nabla |\nabla v|^2|^2 + C_3\left(\int_{\Omega} |\nabla v|^2\right)^2
\\
&\leq \frac{1}{16}\int_{\Omega} |\nabla |\nabla v|^2|^2 + C_4.
\end{align*}
Next, by  the facts $ |\Delta v| \leq \sqrt{2} |D^2 v|$, 
$|\nabla|\nabla v||\leq |D^2 v|$, together with H\"{o}lder inequality and Young's 
inequality, we obtain
\begin{align*}
I_2 + I_3 &=- \int_{\Omega} n |\nabla v|^2 \Delta v
  - \int_{\Omega} n \nabla v \cdot \nabla |\nabla v|^2\\
&\leq C_5\left(\||\nabla v| |D^2 v|\|_{L^2(\Omega)}+\|\nabla|\nabla v|^2\|_{L^2(\Omega)}\right)\|n\|_{L^3(\Omega)}\|\nabla v\|_{L^6(\Omega)}\\
&\leq C_6\left(\||\nabla v| |D^2 v|\|_{L^2(\Omega)}+\|\nabla|\nabla v|^2\|_{L^2(\Omega)}\right)\|n\|_{L^3(\Omega)}\left(\|\nabla|\nabla v|^2\|_{L^2(\Omega)}+1\right)^{\frac{1}{3}}
\\
&\leq \frac{1}{2} \int_{\Omega} |\nabla v|^2 |D^2 v|^2 +\frac{1}{16}\int_{\Omega} |\nabla |\nabla v|^2|^2  + C_7 \int_{\Omega} n^3+C_8,
\end{align*}
where we have used the estimate
\begin{equation}\label{eq3.35}
\begin{aligned}
\int_{\Omega} |\nabla v|^6 &= \left\| |\nabla v|^2 \right\|_{L^3(\Omega)}^3 
\leq C_{9} \left\| \nabla |\nabla v|^2 \right\|_{L^2(\Omega)}^2 \left\| |\nabla v|^2
\right\|_{L^1(\Omega)} + C_{10} \left\| |\nabla v|^2 \right\|_{L^1(\Omega)}^3 \\
&\leq C_{11} \int_{\Omega} |\nabla |\nabla v|^2|^2 + C_{12},
\end{aligned}
\end{equation}
because of \eqref{9.12-1} and \eqref{eq-3.11}. 
Similarly, using \eqref{eq3.6} and \eqref{eq3.35}, we derive $I_4$ and $I_5$ as
\begin{align*}
I_4 &= \int_{\Omega} (u \cdot \nabla v) |\nabla v|^2 \Delta v\\
&\leq \sqrt{2} \int_{\Omega} |u \cdot \nabla v| \, |\nabla v|^2 \, |D^2 v| \\
&\leq C_{13} \||\nabla v| |D^2 v|\|_{L^2(\Omega)}\|\nabla v\|^2_{L^6(\Omega)}\|u\|_{L^6(\Omega)}\\
&\leq C_{14} \||\nabla v| |D^2 v|\|_{L^2(\Omega)}\left(\|\nabla|\nabla v|^2\|_{L^2(\Omega)}+1\right)^{\frac{2}{3}}\\
&\leq \frac{1}{4} \||\nabla v|  |D^2 v|\|_{L^2(\Omega)}^2 + \frac{1}{16}\| \nabla |\nabla v|^2 \|_{L^2(\Omega)}^2 + C_{15},
\end{align*}
and
\begin{align*}
I_5 = \int_{\Omega} (u \cdot \nabla v) \nabla v \cdot \nabla|\nabla v|^2
&\leq C_{16}\|\nabla|\nabla v|^2\|_{L^2(\Omega)}\|\nabla v\|_{L^6(\Omega)}^2\|u\|_{L^6(\Omega)} \\
&\leq \frac{1}{16} \|\nabla|\nabla v|^2\|_{L^2(\Omega)}^2+ C_{17}.
\end{align*}
Substituting the above estimates for $I_1$--$I_5$ into \eqref{(3.45)}, we obtain
\begin{equation}\label{(3.48)}
\frac{d}{dt} \int_{\Omega} |\nabla v|^4
+  \int_{\Omega} |\nabla |\nabla v|^2|^2
+ \int_{\Omega} |\nabla v|^2 |D^2 v|^2
+ 4\int_{\Omega} |\nabla v|^4\leq  C_{18} \int_{\Omega} n^3 +C_{19}.
\end{equation}

Applying the gradient operator $\nabla$ to \eqref{eq1.1}$_2$ and multiplying the 
resulting equations by $4|\nabla c|^2 \nabla c$, we obtain after integrating by parts 
over $\Omega$ that for all $t \in (0, T_{\rm max})$
\begin{equation}\label{eq3.38}
\begin{aligned}
\frac{d}{dt} \int_{\Omega} |\nabla c|^4
&= 4 \int_{\Omega} |\nabla c|^2 \nabla c \cdot \nabla ( \Delta c - n c - u \cdot \nabla c ) \\
&=- 4 \int_{\Omega} |\nabla c|^2 |D^2 c|^2
+ 2 \int_{\Omega} |\nabla c|^2 \Delta |\nabla c|^2 + 4 \int_{\Omega} c n( \nabla |\nabla c|^2 \nabla c + |\nabla c|^2 \Delta c ) \\
& \quad - 4 \int_{\Omega} |\nabla c|^2 \nabla c \cdot \nabla (u \cdot \nabla c) \\
&:= - 4 \int_{\Omega} |\nabla c|^2 |D^2 c|^2 +\sum_{i=1}^{3}J_i,
\end{aligned}
\end{equation}
where we have used the fact 
\begin{gather*}
2\nabla v \cdot \nabla \Delta v = \Delta | \nabla v |^2 - 2 | D^2 v |^2 , \\
-4 \int_{\Omega}
|\nabla c|^2 \nabla c \cdot \nabla(n c) = 4 \int_{\Omega} c n( \nabla |\nabla c|^2 
\nabla c + |\nabla c|^2 \Delta c ).
\end{gather*}
Firstly, using integration by parts and the same calculation procedure as that for $I_1$, 
we obtain
\[
J_1 = 2 \int_{\Omega} |\nabla c|^2 \Delta |\nabla c|^2 
= -2 \int_{\Omega} |\nabla |\nabla c|^2|^2 + 2 \int_{\partial \Omega} |\nabla c|^2
\frac{\partial |\nabla c|^2}{\partial \nu}
 \leq -\frac{7}{4} \int_{\Omega} |\nabla |\nabla c|^2|^2 + C_{20}.
\]
Using that $ |\Delta c| \leq \sqrt{2} |D^2 c|$, H\"{o}lder inequality, 
\eqref{9.12-1}, \eqref{9.12-2}, \eqref{10.20-3}, \eqref{(9.12.3.15)} and 
\eqref{eq-3.11}, we obtain
\begin{align*}
J_2 
&= 4 \int_{\Omega} c n ( \nabla |\nabla c|^2 \cdot \nabla c + |\nabla c|^2 \Delta c )
\\
&\leq 4\|c\|_{L^\infty(\Omega)}\|n\|_{L^3(\Omega)} \|\nabla c\|_{L^6(\Omega)}
( \|\nabla |\nabla c|^2\|_{L^2(\Omega)} + \||\nabla c| \Delta c\|_{L^2(\Omega)} ) \\
&\leq \frac{1}{4} \int_{\Omega} \left(| \nabla |\nabla c|^2|^2
+ |\nabla c|^2 |D^2 c|^2\right) +\varepsilon^{-1/2}\int_{\Omega} n^3
 + C_{21}\varepsilon\int_{\Omega} |\nabla c|^6 + C_{22}.
\end{align*}
Using \eqref{eq-3.11} and applying the same calculation procedure as that for $I_4$ 
and $I_5$, we obtain
\begin{align*}
J_3 
&= - 4 \int_{\Omega} |\nabla c|^2 \nabla c \cdot \nabla (u \cdot \nabla c) \\&=4 \int_{\Omega} (u \cdot \nabla c) |\nabla c|^2 \Delta c
+ 4 \int_{\Omega} (u \cdot \nabla c) \nabla c \cdot \nabla |\nabla c|^2\\
&\leq \int_{\Omega} |\nabla c|^2 \, |D^2 c|^2 + \frac{1}{2} \int_{\Omega} |\nabla |\nabla c|^2|^2 + C_{23}.
\end{align*}
Then, substituting $J_1$--$J_3$ into \eqref{eq3.38}, we obtain
\begin{equation}\label{(3.57)}
\frac{d}{dt} \int_{\Omega} |\nabla c|^4
+ \int_{\Omega} |\nabla |\nabla c|^2|^2
+ \frac{11}{4}\int_{\Omega} |\nabla c|^2 |D^2 c|^2 
\leq \varepsilon^{-1/2}\int_{\Omega} n^3+ C_{24}\varepsilon\int_{\Omega} |\nabla c|^6
+ C_{25}.
\end{equation}

Putting \eqref{(3.48)}, \eqref{eq-3.34}, and \eqref{(3.57)} together leads to
\begin{equation}\label{10.16-1}
\begin{aligned}
&\frac{d}{dt} \int_{\Omega} ( n^2 + |\nabla v|^4 + |\nabla c|^4 )
+ \int_{\Omega} |\nabla n|^2
+  \int_{\Omega} |\nabla |\nabla v|^2|^2 \\
&+ \int_{\Omega} |\nabla v|^2 |D^2 v|^2
+ 4\int_{\Omega} |\nabla v|^4
+  \int_{\Omega} |\nabla |\nabla c|^2|^2
+ \int_{\Omega} |\nabla c|^2 |D^2 c|^2 \\
&\leq (2\varepsilon^{-1/2}+C_{18})\int_{\Omega} n^3
+ C_{26}\varepsilon\int_{\Omega} |\nabla c|^6 + C_{27}\varepsilon\int_{\Omega} 
|\nabla v|^6+ 2\int_{\Omega} nf(n)+C_{28}.
\end{aligned}
\end{equation}
On the other hand, using \eqref{10.20-3}, \eqref{9.12-2} and \eqref{eq-3.11}, for any 
$\varepsilon\in (0,\infty)$ we obtain
\begin{equation}\label{10.16-2}\begin{aligned}
\int_{\Omega} n^3
&\leq \varepsilon \| \nabla n \|_{L^2(\Omega)}^2 \| n\ln n \|_{L^1(\Omega)} 
 + C_\varepsilon \| n \|_{L^1(\Omega)}^3 + C_\varepsilon \\
&\leq \varepsilon C_{29} \| \nabla n \|_{L^2(\Omega)}^2  + C_{30}.
\end{aligned}
\end{equation}
Using \eqref{eq-3.3} and that $ f $ is bounded on finite interval $ [0, \hat{s}] $, 
there exists a constant $ C_{31} $ such that
\[
s^l f(s) \leq
\begin{cases}
-\hat{\mu} \frac{s^{l+2}}{\ln s} \leq 0 & s \geq \hat{s}, \\[3pt]
s^l f(s) \leq C_{31} & s < \hat{s},
\end{cases}
\]
with $l\in \{1,2\}$. This, together with \eqref{eq-3.11}, yields that
\begin{equation}\label{eq-3.42}
\int_{\Omega} n^l f(n) \leq C_{32}.
\end{equation}
Then, substituting \eqref{10.16-2} and \eqref{eq-3.42} ($l=1$) into \eqref{10.16-1} 
and taking $ \varepsilon $ suitable small lead to
\begin{equation}\label{eq-3.43}
\begin{aligned}
&\frac{d}{dt} \int_{\Omega} ( n^2 + |\nabla v|^4 + |\nabla c|^4 )
+ \frac{1}{2}\int_{\Omega} |\nabla n|^2 + \frac{1}{2}\int_{\Omega} |\nabla |\nabla v|^2|^2 \\
&+ \int_{\Omega} |\nabla v|^2 |D^2 v|^2 + 4\int_{\Omega} |\nabla v|^4
+ \frac{1}{2}\int_{\Omega} |\nabla |\nabla c|^2|^2
+ \int_{\Omega} |\nabla c|^2 |D^2 c|^2 \leq C_{33},
\end{aligned}
\end{equation}
where we also used the estimate \eqref{eq3.35} to deal with $\|\nabla v\|_{L^6}$ and 
$\|\nabla c\|_{L^6}$.
Using Poincar\'e inequality, \eqref{10.20-3}, and \eqref{eq-3.11}, we obtain
\begin{gather*}
\|\, |\nabla c|^2 \|_{L^2(\Omega)}^2 
\leq C^2_\mathrm{p} \| \nabla|\nabla c|^2 \|_{L^2(\Omega)}^2 +C_{34}, \\
\| n \|_{L^2(\Omega)}^2 \leq C^2_\mathrm{p} \| \nabla n \|_{L^2(\Omega)}^2 +C_{35}.
\end{gather*}
Substituting the above two inequalities into \eqref{eq-3.43}, we obtain that
\begin{equation}\label{1-18-1}
\begin{aligned}
&\frac{d}{dt} \int_{\Omega} \left( n^2 + |\nabla v|^4 + |\nabla c|^4\right)
\\
&+ \min\left\{1,\frac{1}{4C^2_\mathrm{p}} \right\} \int_{\Omega} \left( n^2
 + |\nabla v|^4 + |\nabla c|^4 + |\nabla v|^2 |D^2 v|^2
 + |\nabla c|^2 |D^2 c|^2\right)\\
&+\frac{1}{4}\int_{\Omega}\left(|\nabla n|^2+|\nabla |\nabla c|^2|^2
 +|\nabla |\nabla v|^2|^2\right)\leq C_{36}.
\end{aligned}
\end{equation}
Applying Lemma \ref{lem2.4} to the above inequality immediately yields the desired 
uniform boundedness estimate \eqref{eq-3.32}.
Using the estimate \eqref{eq-3.32}, and  integrating the inequality \eqref{1-18-1} 
with respect to time  lead to the second desired estimate \eqref{1-18-2} in this lemma.
Now, we complete the proof of this lemma.
\end{proof}

\begin{lemma}\label{lem3.7}
Let $ \alpha \in (\frac{1}{2}, 1) $. Then there exist two positive constants $K_7$ 
and $K_8$ such that
\begin{equation}\label{3.60}
\|\mathcal{A}^\alpha u(\cdot, t) \|_{L^2(\Omega)} \leq K_7 \quad 
\text{for all } t \in (0, T_{\rm max}),
\end{equation}
and consequently
\begin{equation}\label{3.61}
\| u(\cdot, t) \|_{L^\infty(\Omega)} \leq K_8 \quad 
\text{for all } t \in (0, T_{\rm max}).
\end{equation}
\end{lemma}

\begin{proof} 
First of all, we derive the variation-of-constants formula for $u$ from the equation 
\eqref{eq1.1} as
\[
u(x,t) = e^{-t \mathcal{A}} u_0 + \int_{0}^t e^{-(t - s) \mathcal{A}} 
\mathcal{P}(n\nabla \phi(\cdot, s)) ds \quad \text{for all } t \in (0, T_{\rm max}).
\]
Applying the operator $\mathcal{A}^\alpha$ to the above identity and taking the 
$L^2$-norm, one obtains
\begin{equation} \label{3.62}
\| \mathcal{A}^\alpha u(\cdot, t) \|_{L^2(\Omega)}
\leq \| \mathcal{A}^\alpha e^{-t\mathcal{A}} u_0 \|_{L^2(\Omega)}
+ \int_0^t \|\mathcal{A}^{\alpha} e^{-(t - s)\mathcal{A}} \mathcal{P}(n \nabla
\phi(\cdot, s))\|_{L^2(\Omega)} ds.
\end{equation}
Next, we shall use the properties of Stokes operator (see Lemma \ref{lem2.7})
to deal with the two terms on the right hand of \eqref{3.62}. Firstly,
using \eqref{eq-2.9} and the fact $u_0 \in D(\mathcal{A}^\alpha)$, there exists
positive constant $C_1$ such that
\begin{equation}\label{3.63}
\|\mathcal{A}^\alpha e^{-t\mathcal{A}} u_0\|_{L^2(\Omega)}
\leq \| e^{-t\mathcal{A}}\mathcal{A}^\alpha u_0\|_{L^2(\Omega)}
 \leq \|\mathcal{A}^\alpha u_0\|_{L^2(\Omega)} \leq C_1 \quad \text{for
all } t \in (0, T_{\rm max}).
\end{equation}
Since the Helmholtz projection $\mathcal{P}$ is a linear operator from $L^2(\Omega)$
to $L_\sigma^2(\Omega)$, we obtain
\begin{equation}\label{3.64}
\begin{split}
\int_0^t \|\mathcal{A}^{\alpha} e^{-(t - s)\mathcal{A}} \mathcal{P}(n(\cdot, s)
  \nabla \phi)\|_{L^2(\Omega)} ds
&\le C_2  \int_0^t (t - s)^{-\alpha} e^{-\lambda (t - s)}
 \|n\nabla \phi\|_{L^2(\Omega)} ds\\
&\le C_2 \|n\|_{L^2(\Omega)} \|\nabla\phi\|_{L^\infty(\Omega)}
 \int_0^t (t - s)^{-\alpha} e^{-\lambda (t - s)} ds \\
& \le C_3,
\end{split}
\end{equation}
because of the inequality \eqref{eq-2.9}, \eqref{eq-3.32},  the boundedness of
$\nabla\phi$, and $\alpha\in (\frac{1}{2}, 1)$.

Substituting \eqref{3.63} and \eqref{3.64} into \eqref{3.62}, we conclude that there 
exists a positive constant $C_4$ such that
\[
\|\mathcal{A}^\alpha u(\cdot, t)\|_{L^2(\Omega)} \leq C_4 \quad 
\text{for all } t \in (0, T_{\rm max}),
\]
which is exactly the estimate \eqref{3.60}.

Finally, since the domain $D(\mathcal{A}^\alpha)$ is continuously embedded into 
$L^\infty(\Omega)$ for $\frac{1}{2}< \alpha<2$ (see e.g. \cite{Sohr-2001}),
the estimate \eqref{3.60} immediately implies that
\[
\| u(\cdot, t) \|_{L^\infty(\Omega)} \leq C_5 \quad \text{for all } 
t \in (0, T_{\rm max}).
\]
This is \eqref{3.61}.  The proof of Lemma is complete.
\end{proof}

\begin{lemma}\label{lem3.8}
For each $ 2 <q <\infty$, there exists a constant $K_9>0$ such that
\begin{equation}\label{3.66}
\|\nabla c(\cdot, t)\|_{L^q(\Omega)} + \|\nabla v(\cdot, t)\|_{L^q(\Omega)} 
\leq K_9 \quad \text{for all } t \in (0, T_{\rm max}),
\end{equation}
\end{lemma}

\begin{proof} 
We first derive the variation-of-constants formulas for \eqref{eq1.1}$_2$ and 
\eqref{eq1.1}$_3$, respectively, which read
\begin{gather}\label{9.13-1}
c(\cdot, t) = e^{t\Delta} c_0- \int_{0}^t e^{(t - s)\Delta} 
[ u(\cdot, s) \cdot \nabla c(\cdot, s) + n(\cdot, s)c(\cdot, s) ] ds 
\quad \text{for all } t \in (0, T_{\rm max}), \\
\label{9.13-2}
v(\cdot, t) = e^{t\Delta} v_0 - \int_{0}^t e^{(t - s)\Delta} [ u(\cdot, s) 
\cdot \nabla v(\cdot, s) + v(\cdot, s) - n(\cdot, s) ] ds \quad \text{for all
} t \in (0, T_{\rm max}).
\end{gather}
Next, we estimate the $L^q$-norm of $\nabla c$. By applying the gradient operator 
$\nabla$ to \eqref{9.13-1} and taking the $L^q$-norm, together with \eqref{9.12-6} 
and \eqref{1.18-6}, we obtain  for any $2<q <\infty$ that
\begin{align*}
&\|\nabla c(\cdot, t)\|_{L^q(\Omega)} \\&\leq \| \nabla e^{t\Delta} 
 c_0 \|_{L^q(\Omega)} + \int_{0}^t \| \nabla e^{(t - s)\Delta} [ u(\cdot, s) \cdot
\nabla c(\cdot, s) + n(\cdot, s)c(\cdot, s) ] \|_{L^q(\Omega)} ds \\
&\leq C_1\| \nabla c_0\|_{L^q(\Omega)} + C_2 \int_{0}^t (1 + (t - s)^{-1 + \frac{1}{q}})
 e^{-\lambda_1(t - s)} \| u(\cdot, s) \cdot \nabla c(\cdot, s)
\|_{L^2(\Omega)} ds \\
&\quad + C_3 \int_0^t (1 + (t - s)^{-1 + \frac{1}{q}}) e^{-\lambda_2(t - s)} 
 \|  n(\cdot, s)c(\cdot, s)  \|_{L^2(\Omega)} ds\\
&\leq C_4 + C_2 \int_{0}^t (1 + (t - s)^{-1 + \frac{1}{q}}) e^{-\lambda_1(t - s)} 
\| u(\cdot, s) \|_{L^\infty(\Omega)} \|
\nabla c(\cdot, s) \|_{L^2(\Omega)} ds\\
&\quad + C_3 \int_{0}^t (1 + (t - s)^{-1 + \frac{1}{q}}) e^{-\lambda_2(t - s)} 
\| c(\cdot, s) \|_{L^\infty(\Omega)} \| n(\cdot, s) \|_{L^2(\Omega)} ds
\quad \text{for all } t \in (0, T_{\rm max}),
\end{align*}
where $ C_i \ (i = 1, 2, 3, 4) $ and $ \lambda_j \ (j = 1, 2) $ are some positive 
constants. Using \eqref{(9.12.3.15)}, \eqref{eq-3.11},
\eqref{eq-3.32} and \eqref{3.61}, one obtains that
\begin{equation}\label{10.15}
\begin{aligned}
\|\nabla c(\cdot, t)\|_{L^q(\Omega)} 
&\leq C_4 + C_5\int_{0}^t (1 + (t - s)^{-1 + \frac{1}{q}}) e^{-\lambda_1(t - s)} ds\\
&\quad + C_6 \int_{0}^t (1 + (t - s)^{-1 + \frac{1}{q} }) e^{-\lambda_2(t - s)} ds\\
&\leq C_7 \quad \text{for all } t \in (0, T_{\rm max}).
\end{aligned}
\end{equation}
By a similar argument as that applied to the $ c $-equation \eqref{9.13-1}, we deduce 
from \eqref{9.13-2} for all $t \in (0, T_{\rm max})$
\begin{align*}
&\|\nabla v(\cdot, t)\|_{L^q(\Omega)} \\&\leq \| \nabla e^{t\Delta} v_0\|_{L^q(\Omega)} 
 + \int_{0}^t \| \nabla e^{(t - s)\Delta} [ u(\cdot, s) \cdot
\nabla v(\cdot, s) + v(\cdot, s) - n(\cdot, s) ] \|_{L^q(\Omega)} ds \\
&\leq C_8\|\nabla v_0 \|_{L^q(\Omega)} + C_{9} \int_{0}^t
  (1 + (t - s)^{-1 + \frac{1}{q}}) e^{-\lambda_3(t - s)} \| u(\cdot, s) \cdot 
  \nabla v(\cdot,s) \|_{L^2(\Omega)} ds \\
&\quad + C_{10}  \int_{0}^t (1 + (t - s)^{-1 + \frac{1}{q}}) 
 e^{-\lambda_4(t - s)} ( \| v(\cdot, s) \|_{L^2(\Omega)} 
 + \| n(\cdot, s) \|_{L^2(\Omega)}) ds \\
&\leq C_{11} + C_{9} \int_{0}^t (1 + (t - s)^{-1 + \frac{1}{q}}) e^{-\lambda_3(t - s)} 
 \| u(\cdot, s) \|_{L^\infty(\Omega)}\| \nabla v(\cdot, s) \|_{L^2(\Omega)} ds\\
&\quad + C_{10} \int_{0}^t (1 + (t - s)^{-1 + \frac{1}{q}}) e^{-\lambda_4(t - s)}
 ( \| v(\cdot, s) \|_{L^2(\Omega)} + \| n(\cdot, s) \|_{L^2(\Omega)})ds ,
\end{align*}
where $ C_i$ $(i =8, 9, 10, 11) $ and $ \lambda_j \ (j = 3, 4) $ are some positive 
constants. Using \eqref{eq-3.11},
\eqref{eq-3.32} and \eqref{3.61}, one obtains that
\begin{equation}\label{10.15-2}
\begin{aligned}
\|\nabla v(\cdot, t)\|_{L^q(\Omega)}
& \leq C_{11} + C_{12}\int_0^t (1 + (t - s)^{-1 + \frac{1}{q}}) e^{-\lambda_3(t - s)} ds\\
&\quad + C_{13}\int_0^t (1 + (t - s)^{-1 + \frac{1}{q} }) e^{-\lambda_4(t - s)} ds \\
&\leq C_{14} \quad \text{for all } t \in (0, T_{\rm max}).
\end{aligned}
\end{equation}
Then \eqref{3.66}  follows from \eqref{10.15} and \eqref{10.15-2}. 
This completes the proof.
\end{proof}

\begin{lemma}\label{lem3.9}
There exists a constant $K_{10}> 0 $ such that
\begin{equation}\label{9.12-7}
\|n(\cdot,t)\|_{L^3(\Omega)} \leq K_{10} \quad \text{for all } t \in (0, T_{\rm max}).
\end{equation}
\end{lemma}

\begin{proof} 
We multiply both sides of \eqref{eq1.1}$_1$ by $3n^2$ and integrate the resulting 
equation by parts to deduce that
\begin{equation}\label{(3.67)}
\begin{aligned}
\frac{d}{dt} \int_{\Omega} n^3 &= 3 \int_{\Omega} \nabla 
 \cdot  ( \nabla n - \chi n \nabla c + \xi n \nabla v) n^2 
 + 3 \int_{\Omega} n^2 f(n) \\
&\leq -6 \int_{\Omega} n |\nabla n|^2 + 6\chi \int_{\Omega} n^2 \nabla n 
\cdot \nabla c -6\xi \int_{\Omega} n^2 \nabla n \cdot \nabla v + C_1 \\
&\leq -2 \int_{\Omega} n |\nabla n|^2 + 2\int_{\Omega} |n|^4 
 + C_2 \int_{\Omega} |\nabla c|^8 + C_3 \int_{\Omega} |\nabla v|^8 + C_1 \\
&\leq -2 \int_{\Omega} n |\nabla n|^2 + 2\int_{\Omega} |n|^4 
 + C_4 \quad \text{for all } t \in (0, T_{\rm max}),
\end{aligned}
\end{equation}
where we have used Young's inequality, \eqref{3.66}  and \eqref{eq-3.42} 
(taking $l = 2$). Using \eqref{eq-3.32}, we turn to estimate the second term on
the right hand of \eqref{(3.67)} as
\begin{equation}\label{(3.68)}
\begin{aligned}
\| n \|_{L^4(\Omega)}^4
&= \| n^{3/2} \|_{L^{8/3}(\Omega)}^{8/3}\\
&\leq C_5\Big( \| \nabla n^{3/2} \|_{L^2(\Omega)}^{4/3} \|
 n^{3/2} \|_{L^{4/3}(\Omega)}^{4/3}
+ \| n^{3/2} \|_{L^{4/3}(\Omega)}^{8/3}\Big) \\
&\leq \frac{1}{4}\int_{\Omega} n |\nabla n|^2 + C_6.
\end{aligned}
\end{equation}
Substituting \eqref{(3.68)} into \eqref{(3.67)} and using the fact that 
$\|n\|^3_{L^3}\leq \|n\|^4_{L^4}+|\Omega|$, we obtain
\begin{equation}\label{3.67-1}
\begin{aligned}
\frac{d}{dt} \int_{\Omega} n^3 +\int_{\Omega} n^3\leq C_7  \quad \text{for all } 
t \in (0, T_{\rm max}).
\end{aligned}
\end{equation}
Applying Lemma \ref{lem2.4} to \eqref{3.67-1} immediately yields 
\begin{equation}\label{(3.69)}
\begin{aligned}
\| n(\cdot, t) \|_{L^3(\Omega)} \leq C_8 \quad \text{for all } t \in (0, T_{\rm max}).
\end{aligned}
\end{equation}
This completes the proof. 
\end{proof}


\begin{lemma}\label{lem3.10}
There exists a constant $K_{11}> 0$ such that
\begin{equation}\label{(3.70)}
\| c(\cdot, t) \|_{W^{1,\infty}(\Omega)} + \| v(\cdot, t) \|_{W^{1,\infty}(\Omega)}
  \leq K_{11} \quad \text{for all } t \in (t_1, T_{\rm max}),
\end{equation}
where $t_1:=\min\left\{1,\frac{T_{\rm max}}{6}\right\}$.
\end{lemma}

\begin{proof} 
We first derive the variation-of-constants formulas for \eqref{eq1.1}$_2$ and 
\eqref{eq1.1}$_3$, respectively, which read
\begin{gather}\label{9.13-12}
c(\cdot, t) = e^{(t-t_1)\Delta} c(\cdot, t_1)- \int_{t_1}^t
 e^{(t - s)\Delta} [ u(\cdot, s) \cdot \nabla c(\cdot, s) + n(\cdot, s)c(\cdot, s) ] ds 
 \quad \forall t \in (t_1, T_{\rm max}), \\
\label{9.13-22}
v(\cdot, t) = e^{(t-t_1)\Delta} v(\cdot, t_1) - \int_{t_1}^t e^{(t - s)\Delta} 
[ u(\cdot, s) \cdot \nabla v(\cdot, s) + v(\cdot, s) - n(\cdot, s) ] ds \quad 
\forall  t \in (t_1, T_{\rm max}).
\end{gather}

By applying the gradient operator $\nabla$ to \eqref{9.13-12} and taking the 
$L^\infty$-norm, together with and using \eqref{9.12-6}, \eqref{(9.12.3.15)},
 \eqref{3.61}, \eqref{3.66} and \eqref{9.12-7}, we obtain
\begin{equation} \label{3.71}
\begin{aligned}
\|\nabla c(t)\|_{L^\infty(\Omega)}
&\leq \| \nabla e^{(t-t_1)\Delta} c(\cdot,t_1) \|_{L^\infty(\Omega)}
+ \int_{t_1}^t \| \nabla e^{(t - s)\Delta} (u \cdot \nabla c)(s) \|_{L^\infty(\Omega)} ds\\
&\quad + \int_{t_1}^t \| \nabla e^{(t - s)\Delta} (n c)(s)\|_{L^\infty(\Omega)} ds \\
&\leq C_1 ( 1 + t^{-\frac{3}{2}}) \| c(\cdot,t_1)\|_{L^1(\Omega)}
 + C_2 \int_{t_1}^t ( 1 + (t - s)^{-\frac{3}{4}} ) e^{-\lambda_5(t - s)}
  \| (u \cdot \nabla c)(s) \|_{L^4(\Omega)} ds \\
&\quad + C_3 \int_{t_1}^t ( 1 + (t - s)^{-\frac{5}{6}} ) e^{-\lambda_6(t - s)}
 \| (n c)(s) \|_{L^3(\Omega)} ds  \\
&\leq C_4 + C_2 \int_{t_1}^t ( 1 + (t - s)^{-\frac{3}{4}} ) e^{-\lambda_5(t - s)}
 \| u(s) \|_{L^\infty(\Omega)} \| \nabla c(s) \|_{L^4(\Omega)} ds \nonumber\\
&+ C_3 \int_{t_1}^t ( 1 + (t - s)^{-\frac{5}{6}} ) e^{-\lambda_6(t - s)}
 \| n(s) \|_{L^3(\Omega)} \| c(s) \|_{L^\infty(\Omega)} ds \\
&\leq C_5  \quad  \text{for all } t \in (t_1, T_{\rm max}).
\end{aligned}
\end{equation}
By a similar argument as that applied to \eqref{9.13-12} and using Poincar\'e
inequality, we deduce from \eqref{9.13-22} that
\begin{equation} \label{3.72}
\begin{aligned}
&\|\nabla v(t)\|_{L^\infty(\Omega)}\\
&\leq \| \nabla e^{(t-t_1)\Delta} v(\cdot, t_1) \|_{L^\infty(\Omega)}
+ \int_{t_1}^t \| \nabla e^{(t - s)\Delta} (u \cdot \nabla v)(s) \|_{L^\infty(\Omega)}ds
 \\
&\quad + \int_{t_1}^t \| \nabla e^{(t - s)\Delta} v(s) \|_{L^\infty(\Omega)} ds
+ \int_{t_1}^t \| \nabla e^{(t - s)\Delta} n(s) \|_{L^\infty(\Omega)} ds \\
&\leq C_6 ( 1 + t^{-\frac{3}{2}} ) \|v(\cdot, t_1)\|_{L^1(\Omega)}
 + C_7 \int_{t_1}^t ( 1 + (t - s)^{-\frac{3}{4}} ) e^{-\lambda_7(t - s)} \| (u \cdot \nabla
v)(s) \|_{L^4(\Omega)} ds  \\
&\quad + C_8 \int_{t_1}^t ( 1 + (t - s)^{-\frac{5}{6}} ) e^{-\lambda_8(t - s)}
 \| v(s) \|_{L^3(\Omega)} ds \\
&\quad + C_8 \int_{t_1}^t ( 1 + (t - s)^{-\frac{5}{6}} )
e^{-\lambda_8(t - s)} \| n(s) \|_{L^3(\Omega)} ds \\
&\leq C_9 + C_7 \int_{t_1}^t ( 1 + (t - s)^{-\frac{3}{4}} ) e^{-\lambda_7(t - s)}
 \| u(s) \|_{L^\infty(\Omega)} \| \nabla v(s) \|_{L^4(\Omega)} ds  \\
&\quad + C_{10} \int_{t_1}^t ( 1 + (t - s)^{-\frac{5}{6}} ) e^{-\lambda_8(t - s)}
 \| \nabla v(s) \|_{L^3(\Omega)} ds
+ C_{11} \leq C_{12}  \quad \text{for all } t \in (t_1, T_{\rm max}).
\end{aligned}
\end{equation}
Combining \eqref{3.71} and \eqref{3.72}, we obtain
\[
\|\nabla c\|_{L^\infty(\Omega)} + \|\nabla v\|_{L^\infty(\Omega)}
\leq C_{13}   \quad \text{for all } t \in (t_1, T_{\rm max}).
\]
This, together with \eqref{10.20-3} and \eqref{eq-3.2}, leads to \eqref{(3.70)}.
 Now, we complete the proof .
 \end{proof}

\begin{lemma}\label{lem3.11}
There exists a constant $K_{12}>0$ such that
\begin{equation}\label{3.73}
\| n(\cdot, t) \|_{L^\infty(\Omega)} \leq K_{12}
\quad \text{for all } t \in (0, T_{\rm max}).
\end{equation}
\end{lemma}

\begin{proof} By the fact that \( \nabla \cdot u \equiv 0 \), we derive the 
variation-of-constants formula for $n$ from \eqref{eq1.1}$_1$. 
For each \( t \in (0,T_{\rm max}) \), this formula reads
\begin{equation}\label{3.74-1}
\begin{aligned}
n(\cdot, t) 
&= e^{(t-t_0)\Delta} n(\cdot,t_0) -\chi \int_{t_0}^t  e^{(t - s)\Delta} \nabla 
 \cdot (n(\cdot, s) \nabla c(\cdot, s)) ds \\
&\quad +\xi \int_{t_0}^t e^{(t - s)\Delta} \nabla \cdot (n(\cdot, s) 
 \nabla v(\cdot, s)) ds \\
&\quad - \int_{t_1}^t  e^{(t - s)\Delta} \nabla \cdot (n(\cdot, s) u(\cdot, s))  ds
+ \int_{t_0}^te^{(t - s)\Delta} f(n(\cdot, s))ds\\
&:= n_1(\cdot,t) + n_2(\cdot,t) + n_3(\cdot,t)+ n_4(\cdot,t)+ n_5(\cdot,t),
\end{aligned}
\end{equation}
where $t_0 := (t-1)^+$. Next, we proceed to estimate each of the five terms on the 
right-hand side of \eqref{3.74-1} one by one. First, by the maximum principle for 
the Neumann heat semigroup, we can estimate
\[
\|n_1(\cdot,t)\|_{L^\infty(\Omega)} \leq \|n_0\|_{L^\infty(\Omega)} \quad 
\text{for\ all }\ t \in (0,1],
\]
whereas if $t > 1$, then standard $L^p \to L^q$ estimates for the Neumann heat 
semigroup provide $C_1> 0$ such that
\[
\|n_1(\cdot,t)\|_{L^\infty(\Omega)} \leq C_1 (t - t_0)^{-1} \|n(\cdot,t_0)
\|_{L^1(\Omega)}\leq C_1 \|n\|_{L^1(\Omega)} \leq C_1M \quad 
\text{for\ all }\ t \in (1,T_{\rm max}],
\]
holds because of \eqref{10.20-3}. Thus, we obtain for any $t\in (0,T_{\rm max})$
\begin{equation}\label{3.75}
\|n_1(\cdot,t)\|_{L^\infty(\Omega)} \leq C_{2}.
\end{equation}
Then, using \eqref{3.61}, \eqref{9.12-7}, \eqref{(3.70)} and H\"{o}lder inequality, 
one applies the known smoothing properties of the semigroup
\eqref{9.12-6-2} to obtain for any $t\in (0,T_{\rm max})$ that
\begin{equation}
\begin{aligned}\label{3.76}
\| n_2(\cdot,t) \|_{L^\infty(\Omega)}
&\leq C_{3} \int_{t_0}^t ( 1 + (t - s)^{-\frac{5}{6}} ) 
 e^{-\lambda_9 (t-s)} \| n(\cdot, s) \nabla c(\cdot, s) \|_{L^{3}(\Omega)} ds \\
&\leq C_{3} \int_{t_0}^t ( 1 + (t - s)^{-\frac{5}{6}} ) 
 e^{-\lambda_9 (t-s)} \| n(\cdot, s) \|_{L^3(\Omega)} \| \nabla c(\cdot, s) 
 \|_{L^\infty(\Omega)} ds \leq C_{4},
\end{aligned}
\end{equation}
\begin{equation}
\begin{aligned}\label{3.77}
\| n_3(\cdot,t) \|_{L^\infty(\Omega)}
&\leq C_{5} \int_{t_0}^t ( 1 + (t - s)^{-\frac{5}{6}} ) e^{-\lambda_{10} (t-s)} \| n(\cdot, s) \nabla v(\cdot, s) \|_{L^{3}(\Omega)} ds \\
&\leq C_{5} \int_{t_0}^t ( 1 + (t - s)^{-\frac{5}{6}} ) e^{-\lambda_{10} (t-s)} \| n(\cdot, s) \|_{L^3(\Omega)} \| \nabla v(\cdot, s)
\|_{L^\infty(\Omega)} ds \leq C_{6},
\end{aligned}
\end{equation}
and
\begin{equation}
\begin{aligned}\label{3.78}
\| n_4(\cdot,t) \|_{L^\infty(\Omega)}
&\leq C_{7} \int_{t_0}^t ( 1 + (t - s)^{-\frac{5}{6}} ) e^{-\lambda_{11} (t - s)} \| n(\cdot, s) u(\cdot, s) \|_{L^3(\Omega)} ds \\
&\leq C_{7} \int_{t_0}^t ( 1 + (t - s)^{-\frac{5}{6}} ) e^{-\lambda_{11} (t - s)} \| n(\cdot, s) \|_{L^3(\Omega)} \| u(\cdot, s) \|_{L^\infty(\Omega)} ds
\leq C_{8},
\end{aligned}
\end{equation}
where $\lambda_j > 0  \ (j = 9, 10, 11) $ denotes the first nonzero eigenvalue of
 $-\Delta$ under homogeneous Neumann boundary condition.
Finally, by virtue of the assumptions on the sub-logistic term $f(s)$, 
it is easy to check that there exists a constant $C_9$ suchu $f(s) \leq C_9 $ for 
all \( s \in \mathbb{R}^+ \). Using the maximum principle once more, we have for any 
$t\in (0,T_{\rm max})$
\begin{equation}\label{3.79}
n_5(\cdot, t) \leq \int_{t_0}^t e^{(t-s)\Delta} C_9 ds = C_9(t - t_0) \leq C_9.
\end{equation}
Combination with \eqref{3.75}--\eqref{3.79}, we obtain for any $t\in (0,T_{\rm max})$
\[
\|n(\cdot, t)\|_{L^\infty(\Omega)} = \sup_{x\in\Omega} n(x,t) 
\leq \sum^5_{i=1}\sup_{x\in\Omega} n_i(x,t)\leq C_{10}.
\]
This is the desired estimate \eqref{3.73}. The proof is complete.
\end{proof}

With the above estimates in hand, we are in a position to show Theorem \ref{Th1.1}.

\begin{proof}[Proof of Theorem \ref{Th1.1}]
Combining \eqref{3.60}, \eqref{(3.70)} and \eqref{3.73}, we can obtain the time 
independent boundedness for 
$\|n(\cdot, t)\|_{L^\infty(\Omega)} + \|c(\cdot, t)\|_{W^{1,\infty}(\Omega)} 
+ \|v(\cdot, t)\|_{W^{1,\infty}(\Omega)}+ \|A^\alpha u (\cdot, t)\|_{L^2(\Omega)}$. 
By the extendibility alternative in Lemma \ref{lm2.8}, we thus infer that $T^*=\infty$, 
i.e., the solution $(n, c, v, u, P)$ is global in time. Furthermore, in view of 
\eqref{3.61}, \eqref{(3.70)} and \eqref{3.73},  we directly verify that the boundedness 
assertion in \eqref{eq-1.12} holds true. This completes the proof.
\end{proof}

\subsection*{Acknowledgements}
The authors are  grateful to the anonymous
reviewers for their carefully reading and valuable suggestions which greatly 
improved this article. This work is partially supported by the
Shandong Provincial Natural Science Foundation
 (No. ZR2022JQ06) as well as by the National Natural
Science Foundation of China (No. 11601215, No. 12226347), and 
by the Natural Science Foundation of Liaoning Province (No. 2023-MS-137).

\begin{thebibliography}{10}

\bibitem{ma-2} K. Baghaeia, A. Khelghatib; 
\emph{Global existence and boundedness of classical solutions for a chemotaxis model with 
consumption of chemoattractant and logistic source},
 Math. Models Methods Appl. Sci., 40 (2017), 3799-3807.

\bibitem{Cao2015} X. Cao; 
\emph{Global bounded solutions of the higher-dimensional Keller-Segel system under 
smallness conditions in optimal spaces}, 
Discrete Contin. Dyn. Syst. S., 35 (2015), 1891-1904.

\bibitem{Dolbeault-2004} J. Dolbeault, B. Perthame;
\emph{Optimal critical mass in the two-dimensional Keller-Segel model in $\mathrm{R}^2$}, 
C. R. Math. Acad. Sci. Paris, 339 (2004), 611-616.

\bibitem{G-8} R. Duan, A. Lorz, P. Markowich;
\emph{Global solutions to the coupled chemotaxis-fluid equations}, 
Comm. Partial Differential Equations, 35 (9) (2010), 1635-1673.

\bibitem{Espejo} E. Espejo, T. Suzuki;
\emph{Reaction terms avoiding aggregation in slow fluids}, 
Nonlinear Anal. Real World Appl., 21 (2015), 110-126.

\bibitem{G-3} H. Gajewski, K. Zacharias;
\emph{Global behavior of a reaction-diffusion system modelling chemotaxis}, 
Math. Nachr., 195 (1998), 77-114.

\bibitem{Giga-1991} Y. Giga, H. Sohr;
\emph{Abstract $L^p$ estimates for the Cauchy problem with applications to the 
Navier-Stokes equations in exterior domains}, J. Funct. Anal., 102 (1991), 72-94.

\bibitem{Herrero-1997} M. Herrero, J. Vel\'{a}zquez;
\emph{Finite-time aggregation into a single point in a reaction-diffusion system}, 
Nonlinearity, 10 (1997), 1739-1754.

\bibitem{G-9} M. Herrero, J. Vel\'{a}zquez;
\emph{A blow-up mechanism for a chemotaxis model}, 
Ann. Sc. Norm. Super. Pisa Cl. Sci., 24 (1997), 633-683.

\bibitem{Hillen} T. Hillen, K. Painter;
\emph{A user's guide to PDE models for chemotaxis}, 
J. Math. Biol., 58 (2009), 183-217.

\bibitem{G-5} D. Horstmann, G. Wang;
\emph{Blow-up in a chemotaxis model without symmetry assumptions}, 
European J. Appl. Math., 12 (2001), 159-177.

\bibitem{Horstmann2005} D. Horstmann, M. Winkler;
\emph{Boundedness vs. blow-up in a chemotaxis system}, 
J. Differential Equations, 215 (2005), 52-107.

\bibitem{Jin} C. Jin;
\emph{Large time periodic solutions to coupled chemotaxis-fluid models}, 
Z. Angew. Math. Phys., 68 (6) (2017), 137.

\bibitem{Xiang2018-C} H. Jin, T. Xiang;
\emph{Chemotaxis effect vs. logistic damping on boundedness in the 2-D minimal 
Keller-Segel model}, C. R. Math. Acad. Sci. Paris, 356 (2018), 875-885.

\bibitem{Keller-1970} E. Keller, L. Segel;
\emph{Initiation of slime mold aggregation viewed as an instability}, 
J. Theor. Biol., 26 (1970), 399-415.

\bibitem{Keller-1970-2} E. Keller, L. Segel;
\emph{Model for Chemotaxis}, J. Theor. Biol., 30 (1970), 225-234.

\bibitem{kozono} H. Kozono, M. Miura, Y. Sugiyama;
\emph{Existence and uniqueness theorem on mild solutions to the Keller-Segel 
system coupled with the Navier-Stokes fluid}, 
J. Funct. Anal., 270 (5) (2016), 1663-1683.

\bibitem{ma-13} J. Lankeit;
\emph{Long-term behaviour in a chemotaxis-fluid system with logistic source}, 
Math. Models Methods Appl. Sci., 26 (2016), 2071-2109.

\bibitem{ma-14} J. Lankeit, Y. Wang;
\emph{Global existence, boundedness and stabilization in a high-dimensional 
chemotaxis system with consumption}, Discrete Contin. Dyn. Syst., 37 (2017), 6099-6121.

\bibitem{s1468} X. Li, Y. Xiao;
\emph{Global existence and boundedness in a 2D Keller-Segel-Stokes system}, 
Nonlinear Anal. Real World Appl., 37 (2017), 14-30.

\bibitem{s1468-17} J. Liu, A. Lorz;
\emph{A coupled chemotaxis-fluid model: Global existence},
 Ann. Inst. Henri Poincar\'e Nonlinear Anal., 28 (5) (2011), 643-652.

\bibitem{Luo} W. Luo, Z. Li;
\emph{Global boundedness in a two-dimensional chemotaxis-Navier-Stokes system with 
double chemical signals and nonlinear diffusion}, 
Nonlinear Anal. Real World Appl., 87 (2026), 104415.

\bibitem{ma} Q. Ma;
\emph{Sub-logistic source can prevent blow-up in the chemotaxis-Navier-Stokes system}, 
Discrete Contin. Dyn. Syst. B., 7 (2025), 30.

\bibitem{Nagai} T. Nagai;
\emph{Blow-up of radially symmetric solutions of a chemotaxis system}, 
Adv. Math. Sci. Appl., 5 (1995), 581-601.

\bibitem{G-4} T. Nagai, T. Senba, K. Yoshida;
\emph{Application of the Trudinger-Moser inequality to a parabolic system of chemotaxis},
Funkc. Ekvac. Ser. Int., 40 (1997), 411-433.

\bibitem{Nirenberg-1959} L. Nirenberg;
\emph{On elliptic partial differential equations},
 Ann. Scuola Norm. Sup. Pisa, 13 (3) (1959), 115-162.

\bibitem{Osaki-2001} K. Osaki, A. Yagi;
\emph{Finite dimensional attractor for one-dimensional Keller-Segel equations}, 
Funkcial. Ekvac., 44 (2001), 441-469.

\bibitem{Osaki-2002} K. Osaki, T. Tsujikaw, A. Yagi, M. Mimura;
\emph{Exponential attractor for a chemotaxis growth system of equations}, 
Nonlinear Anal., 51 (2002), 119-144.

\bibitem{Painter} K. Painter, P. Maini, H. Othmer;
\emph{Development and applications of a model for cellular response to multiple
 chemotactic cues}, J. Math. Biol., 41 (4) (2000), 285-314.

\bibitem{G-35} G. Ren, B. Liu;
\emph{A new result for global solvability to a two-dimensional attraction-repulsion
 Navier-Stokes system with consumption of chemoattractant}, 
J. Differential Equations, 336 (2022), 126-166.

\bibitem{G-6} T. Senba, T. Suzuki;
\emph{Parabolic system of chemotaxis: blow up in a finite and the infinite time},
 Methods Appl. Anal., 8 (2001), 349-367.

\bibitem{Sohr-2001} H. Sohr;
\emph{The Navier-Stokes Equations: An Elementary Functional Analytic Approach}, 
Birkh\"auser Advanced Texts Basler Lehrb\"{u}cher, (2001).

\bibitem{ma-19} C. Stinner, C. Surulescu, M. Winkler;
\emph{Global weak solutions in a PDE-ODE system modeling multiscale cancer cell 
invasion}, SIAM J. Math. Anal., 46 (2014), 1969-2007.

\bibitem{Tuval-2005} I. Tuval, L. Cisneros, C. Dombrowski, C. Wolgemuth, J. Kessler, R. Goldstein; Bacterial swimming and oxygen transport near contact lines, Proc. Natl. Acad. Sci. U.S.A., 102 (2005), 2277-2282.

\bibitem{Winkler2014} Y. Tao, M. Winkler;
\emph{Energy-type estimates and global solvability in a two-dimensional 
chemotaxis-haptotaxis model with remodeling of non-diffusible attractant}, 
J. Differential Equations, 257 (2014), 784-815.

\bibitem{s1468-30} Y. Tao, M. Winkler;
\emph{Blow-up prevention by quadratic degradation in a two-dimensional 
Keller-Segel-Navier-Stokes system}, Z. Angew. Math. Phys., 67 (2016), 138.

\bibitem{s1468-31} Y. Tao, M. Winkler;
\emph{Boundedness and decay enforced by quadratic degradation in a three-dimensional 
chemotaxis-fluid system}, Z. Angew. Math. Phys., 66 (5) (2015), 2555-2573.

\bibitem{Winkler2007} J. Tello, M. Winkler;
\emph{A chemotaxis system with logistic source}, 
Comm. Partial Differential Equations, 32 (2007), 849-877.

\bibitem{ma-24} L. Wang, C. Mu, X. Hu, P. Zheng;
\emph{Boundedness in a quasilinear chemotaxis model with consumption of 
chemoattractant and logistic source}, Appl. Anal., 97 (2018), 756-774.

\bibitem{Wang-Xiang} Y. Wang, Z. Xiang;
\emph{Global existence and boundedness in a Keller-Segel-Stokes system involving 
a tensor-valued sensitivity with saturation}, 
J. Differential Equations, 259 (12) (2015), 7578-7609.

\bibitem{Winkler-2015} M. Winkler;
\emph{Boundedness and large time behavior in a three-dimensional chemotaxis-Stokes 
system with nonlinear diffusion and general sensitivity}, 
Calc. Var. Partial Differential Equations, 54 (2015), 3789-3828.

\bibitem{G-10} M. Winkler;
\emph{Aggregation vs. global diffusive behavior in the higher-dimensional Keller-Segel
 model}, J. Differential Equations, 248 (2010), 2889-2905.

\bibitem{W2020-28} M. Winkler;
\emph{Global large-data solutions in a chemotaxis-(Navier-)Stokes system modeling 
cellular swimming in fluid drops}, 
Commun. Partial Differential Equations, 37 (2012), 319-351.

\bibitem{W2020-29} M. Winkler;
\emph{Global weak solutions in a three-dimensional chemotaxis-Navier-Stokes system}, 
Ann. Inst. Henri. Poincar\'e Anal. Non Lin\'eaire., 33 (2016), 1329-1352.

\bibitem{W2020-30} M. Winkler;
\emph{How far do chemotaxis-driven forces influence regularity in the Navier-Stokes 
system?} Trans. Am. Math. Soc., 369 (2017), 3067-3125.

\bibitem{W2020-31} M. Winkler;
\emph{Stabilization in a two-dimensional chemotaxis-Navier-Stokes system}, 
Arch. Ration. Mech. Anal., 211 (2014), 455-487.

\bibitem{Winkler2010} M. Winkler; Boundedness in the higher-dimensional parabolic-parabolic chemotaxis system with logistic source, Comm. Partial Differential Equations, 35 (2010), 1516-1537.

\bibitem{Winkler2023} M. Winkler;
\emph{$L^1$ solutions to parabolic Keller-Segel systems involving arbitrary superlinear 
degradation}, Ann. Sc. Norm. Super. Pisa Cl. Sci., 24 (2023), 141-172.

\bibitem{Winkler2019} M. Winkler;
\emph{A three-dimensional Keller-Segel-Navier-Stokes system with logistic source: 
Global weak solutions and asymptotic stabilization}, J. Funct. Anal., 5 (2019), 2761339.

\bibitem{Xiang2018-2} T. Xiang;
\emph{Sub-logistic source can prevent blow-up in the 2D minimal Keller-Segel 
chemotaxis system}, J. Math. Phys., 59 (2018), 081502.

\bibitem{Xiang2018-J} T. Xiang;
\emph{How strong a logistic damping can prevent blow-up for the minimal Keller-Segel 
chemotaxis system?} J. Math. Anal. Appl., 459 (2018), 1172-1200.

\bibitem{Global-1} L. Xie, Y. Xu;
\emph{Global existence and stabilization in a two-dimensional chemotaxis-Navier-Stokes 
system with consumption and production of chemosignals}, 
J. Differential Equations, 354 (2023), 325-372.

\bibitem{Zhang} Q. Zhang, Y. Zhang;
\emph{On the global well-posedness for the 2D incompressible Keller-Segel-Navier-Stokes 
equations}, Z. Angew. Math. Mech., 11 (2019), 99.

\end{thebibliography}

\end{document}

